\documentclass{article}
\usepackage[utf8]{inputenc}
% %\input{GrandMacros.tex}
% %\input{Macro.tex}
\usepackage{amsthm}
\usepackage{xcolor}
\newtheorem{remark}{Remark}
\newtheorem{lemma}{Lemma}
\newtheorem{theorem}{Theorem}
\newtheorem{proposition}{Proposition}
\newtheorem{corollary}{Corollary}
\usepackage{graphicx,verbatim,array,multicol, subfigure, color, lscape, mathrsfs, dsfont,url}
%\usepackage{psfrag, amsmath, amsfonts, epsfig, fancybox, setspace,soul, cancel,biometri}
\usepackage{psfrag, amsmath, amsfonts, epsfig, fancybox, setspace, soul, amsthm, longtable, threeparttable, threeparttablex, makecell, xcolor, pdflscape, tabu, booktabs, natbib, amssymb}
\usepackage[normalem]{ulem}
\usepackage{cancel}
\usepackage[utf8]{inputenc}
\parindent 16pt
\include{common-defs}
\newcommand{\PP}{\mathbb{P}}

\input{typography-preamble}
% \def\cor{\text{cor}}
% \def\cF{\cF}

\title{Surrogate inference via random probability distributions: evaluation of surrogate quality in future studies}


\begin{document}
\maketitle

\section{Introduction}

Surrogate markers promise to accelerate clinical research by replacing expensive or long-term outcomes with earlier-measured alternatives. A central question is whether a surrogate validated in one study will perform well in future trials with different populations. This is fundamentally a question of \textbf{transportability}: will the surrogate-outcome relationship observed in the current study hold when the population, treatment effect heterogeneity, or covariate distributions change?

Traditional validation methods---proportion of treatment effect (PTE; \cite{chen2003pte}), within-study correlation, principal stratification \cite{frangakis2002principal}, and causal mediation \cite{vanderweele2015mediation}---are designed for this prospective use case. However, these methods \textbf{assume} that observed relationships will transport to future settings. PTE assumes the proportion of treatment effect remains stable; mediation assumes pathway structures persist; principal stratification assumes stratum definitions transfer. When these assumptions are violated---as may occur with covariate shift, treatment effect heterogeneity, or selection mechanisms---inference can be misleading. \cite{parast2024} identified ``limited work on transportability of surrogate knowledge from one study to another'' as a critical methodological gap.

We introduce a framework for \textbf{local geometric evaluation} of transportability. [TO BE UPDATED: Describe framework as characterizing surrogate quality functionals $\phi(\cF)$ across distributions $\Q \sim \cF$ restricted to local geometries $\cU(\PP_0, \lambda; d)$ for any distance metric $d$. The ``local'' constraint (distance $\leq \lambda$) focuses on plausible shifts; the ``geometric'' aspect (choice of $d$) captures different shift mechanisms; ``evaluation'' means explicit characterization of how surrogate quality varies without assuming transportability holds.]


\section{Setting}\label{sec:setting}
Suppose we observe $n$ iid realizations of $\bO = (\bX, A, S, Y) \sim \PP_0$ in the current study, where $A \in \{0, 1\}$ is a binary treatment, $Y \in \cY \subset \R$ is the outcome of interest, $S \in \cS \subset \R$ is a surrogate marker that may replace $Y$ in a future study, and $\bX \in \cX \subset \R^p$ are pre-treatment covariates ($\bX$ may be omitted in a randomized trial). The distribution $\PP_0$ is the law of the observables. We measure treatment effects on the surrogate and the outcome via
\begin{align*}
    &\Delta_Y(\PP_0) = \E_{\PP_0}\{Y(1) - Y(0)\}\\
    &\Delta_S(\PP_0) = \E_{\PP_0}\{S(1) - S(0)\}
\end{align*}

Under suitable identification assumptions, $\Delta_S(\PP_0)$ and $\Delta_Y(\PP_0)$ are identified and can be estimated by standard methods (e.g., difference in means or influence-function-based estimators). However, for the purposes of determining surrogate quality, we wish to know how the surrogate would perform in \emph{future} studies. Consider a future study with distribution $\Q$; we would observe (an estimate of) $\Delta^S(\Q) = \E_{\Q}\{S(1) - S(0)\}$ but not $Y$, and we wish to know what $\Delta_S(\Q)$ implies about the unobservable $\Delta_Y(\Q) = \E_{\Q}\{Y(1) - Y(0)\}$. 

To formalize this, we must begin to characterize a class of future studies to which we might wish to generalize. We treat the study-specific distribution as a draw from a random probability measure $\cF$, so that both the current distribution $\PP_0$ and any future study distribution $\Q$ are (conceptually) draws from $\cF$. Our goal is to make inference about functionals $\phi(\cF)$ of the distribution of studies using the observed data from the current study. An example of a functional that may be of interest is the correlation between the two treatment effects $\Phi(\cF) = \frac{\Cov_{\cF}\{\Delta_S(\Q), \Delta_Y(\Q)\}}{\sqrt{\Var_\cF\{\Delta_S(\Q)\}\Var_\cF\{\Delta_Y(\Q)\}}}$. Our goal is to estimate $\Phi(\cF)$ for an appropriately chosen $\cF$, thereby characterizing how the surrogate--outcome relationship may behave in future studies.

\subsection{Local geometries for future studies}

Let $\Omega = \cX \times \{0,1\} \times \cS \times \cY$ denote the observed data space, and let $\cM_1(\Omega)$ denote the space of all probability measures on $\Omega$. The current study distribution $\PP_0$ is one element of $\cM_1(\Omega)$; future study distributions $\Q$ are also elements of this space.

We operate under the assumption that future studies are likely to bear \emph{some} resemblance to the current study -- if the current study is to be informative of the future study, they should have some similarity. We encode this notion of similarity through a \emph{distance metric} $d: \cM_1(\Omega) \times \cM_1(\Omega) \to [0,\infty)$ that quantifies how different two study distributions are. Given a distance metric $d$ and a tolerance level $\lambda \geq 0$, we define the \emph{local geometry} around $\PP_0$ as
\begin{align*}
    U(\PP_0, \lambda; d) = \{\Q \in \cM_1(\Omega) : d(\Q, \PP_0) \leq \lambda\}.
\end{align*}
This is the set of all study distributions within distance $\lambda$ of the current study. The parameter $\lambda$ controls how far future studies can deviate from $\PP_0$: when $\lambda = 0$, only $\Q = \PP_0$ is admissible; as $\lambda$ increases, more disparate future studies are considered plausible.

To characterize surrogate quality across a range of plausible future studies, we treat the future study distribution $\Q$ as random, drawn from some probability measure $\mu$ on the local geometry $U(\PP_0, \lambda; d)$. We write
\begin{align*}
    \Q \sim \mu \text{ on } U(\PP_0, \lambda; d)
\end{align*}
to denote that $\Q$ is a random draw from $\mu$ with support on $U(\PP_0, \lambda; d)$. A natural choice is the \emph{uniform measure} on $U(\PP_0, \lambda; d)$, which treats all study distributions at the same distance from $\PP_0$ as equally plausible. This uniform measure can be approximated via Markov chain Monte Carlo methods (e.g., hit-and-run sampling) when the geometry is a convex body.

The choice of $\mu$ determines which directions of deviation from $\PP_0$ are probabilistically emphasized. A uniform measure provides a conservative, non-informative assessment: surrogate quality is averaged over all directions of perturbation equally. Alternative choices (e.g., measures that concentrate near $\PP_0$) may be appropriate if prior knowledge suggests certain types of study differences are more plausible than others.

% \subsubsection{Choosing the innovation distribution $\mu$}
%
% The interpretation of surrogate quality depends critically on the choice of $\mu$, not just $\lambda$. If $\mu$ concentrates mass near $\PP_0$ (e.g., $\mu = \text{Dirichlet}(\alpha \bp_0)$ with large $\alpha$), then finding $\phi(\cF_\lambda) \geq c$ for large $\lambda$ provides only weak evidence---it merely says the surrogate performs well for future studies that closely resemble $\PP_0$. A more stringent test uses a non-informative $\mu$ that treats all innovations $\tilde{\PP} \in \cM_1(\Omega)$ equally, providing no favored direction of deviation from $\PP_0$.
%
% For computational tractability and to define a natural non-informative prior, we assume finite support:

\noindent\textbf{Assumption (Finite support).}\label{ass:finite}
The outcome space $\Omega$ has finite support. Equivalently, the observed data naturally partitions into a finite number of cells (e.g., based on observed covariate values, treatment, surrogate, and outcome levels).

This assumption is reasonable in practice: with $n$ observations, empirical distributions are inherently discrete; many applications involve categorical or discretized outcomes; and finite support makes inference tractable. Under finite support, $\cM_1(\Omega)$ is a $(k-1)$-simplex (where $k = |\Omega|$).

\subsubsection{Treatment effect heterogeneity as the fundamental object}

While the finite support assumption allows for tractable computation, it is important to recognize that the \emph{fundamental source of variation} across studies is heterogeneity in treatment effects, not arbitrary partitions of the data space. For continuous covariates $\bX$, treatment effects $\tau_S(\bX) = S(1) - S(0)$ and $\tau_Y(\bX) = Y(1) - Y(0)$ are continuous functions. Future studies differ from the current study primarily through different distributions of $\bX$, which induce different distributions of treatment effects.

The finite support assumption operationalizes this by discretizing $\bX$ into $J$ bins and treating each bin as a cell. The choice of $J$ represents an approximation: as $J \to \infty$ (with $J = o(n)$), the discretized model approaches the continuous treatment effect function. In practice, $J$ should be chosen large enough to capture meaningful heterogeneity but small enough to avoid over-discretization (which increases estimation noise).

This perspective clarifies that there are no ``true'' types---only varying treatment effects. Discretization schemes that align with treatment effect heterogeneity (e.g., using random forests to partition based on estimated $\tau(\bX)$) will provide better approximations than arbitrary covariate-based discretizations. Different discretization schemes explore different aspects of the treatment effect distribution; using multiple schemes and taking their minimum (ensemble approach) provides a more comprehensive assessment of worst-case surrogate performance.

% and we can define a canonical non-informative prior.
%
% Under finite support, the natural choice for a non-informative $\mu$ is the uniform Dirichlet:
% \begin{align}
%     \mu = \text{Dirichlet}(1, \ldots, 1).
% \end{align}
% This distribution is uniform on the simplex $\cM_1(\Omega)$, treating all probability vectors equally. It is the maximum entropy prior given no constraints, providing the most conservative (stringent) assessment of surrogate quality. When $\phi(\cF_\lambda) \geq c$ under this choice of $\mu$, the surrogate performs well across a broad, evenly-distributed set of future studies within distance $\lambda$ of $\PP_0$.
%
% \subsubsection{How similar must a future study be for the surrogate to remain good?}
%
% A central question is: \emph{how different can a future study be from the current study while the surrogate remains good?} We formalize this by asking for what values of $\lambda$ the functional $\phi(\cF_\lambda)$ stays above a desired threshold. For example, if we define ``good surrogate'' as correlation of treatment effects above 0.8, we want to find the largest $\lambda$ for which $\phi(\cF_\lambda) \geq 0.8$. This $\lambda$ tells us how far a future study can deviate from the current study (in total variation distance) while maintaining adequate surrogate quality.
%
% To make this operational, we need to understand how functionals $\phi(\cF_\lambda)$ behave as $\lambda$ varies. The key insight is that the treatment effects $\Delta_S(\Q)$ and $\Delta_Y(\Q)$ depend continuously on $\Q$ (in the sense that small changes in $\Q$ produce small changes in the treatment effects), and thus functionals of the treatment effects also vary smoothly with $\lambda$. The following assumptions and result formalize this.
%
% \noindent\textbf{Assumption 3 (Bounded outcomes).}
% $S \in [s_{\min}, s_{\max}]$ and $Y \in [y_{\min}, y_{\max}]$ are bounded.
%
% Under this assumption, treatment effects are Lipschitz continuous in $\Q$. For any $\Q, \Q'$,
% \begin{align}\label{eq:lipschitz}
%     |\Delta_S(\Q) - \Delta_S(\Q')| &\leq 2(s_{\max} - s_{\min}) \, d_{\mathrm{TV}}(\Q, \Q')\\
%     |\Delta_Y(\Q) - \Delta_Y(\Q')| &\leq 2(y_{\max} - y_{\min}) \, d_{\mathrm{TV}}(\Q, \Q'). \nonumber
% \end{align}
% This follows because expectations of bounded functions are Lipschitz in total variation: $|\E_{\Q} f - \E_{\Q'} f| \leq (\sup f - \inf f) \, d_{\mathrm{TV}}(\Q, \Q')$. Under SUTVA and consistency, $S(1) - S(0)$ ranges over $[-(s_{\max} - s_{\min}), (s_{\max} - s_{\min})]$, yielding the constant $2(s_{\max} - s_{\min})$, and similarly for $Y$.
%
% \begin{theorem}[Surrogate quality under perturbations]\label{thm:epsilon}
% Suppose ``good surrogate'' is defined by $\phi(\cF_\lambda) \geq c$ for some threshold $c$ (e.g., $\phi$ is the correlation of treatment effects and $c = 0.8$). Under Assumption 3, the functional $\phi(\cF_\lambda)$ varies continuously with $\lambda$. If $\phi(\cF_\lambda) \geq c$ for some $\lambda$, then any future study $\Q$ with $d_{\mathrm{TV}}(\Q, \PP_0) \leq \lambda$ will have $|\phi(\Q) - \phi(\cF_\lambda)| \leq L \lambda$ for some constant $L$ depending on $\phi$ and regularity conditions (e.g., non-degenerate variances for correlation).
% \end{theorem}
% \begin{proof}
% By \eqref{eq:lipschitz}, the treatment effects $(\Delta_S(\Q), \Delta_Y(\Q))$ are Lipschitz continuous in $\Q$ with respect to $d_{\mathrm{TV}}$. The functionals $\phi$ of interest (correlation, probability of concordance, conditional expectation) are continuous---and under bounded $S$, $Y$ and non-degenerate variances, Lipschitz---as functions of the bivariate distribution of treatment effects. Since $\Q \in \cF_\lambda(\PP_0)$ implies $d_{\mathrm{TV}}(\Q, \PP_0) \leq \lambda$ by \eqref{eq:distance-bound}, the distribution of $(\Delta_S(\Q), \Delta_Y(\Q))$ varies continuously with $\lambda$, and composing the Lipschitz maps yields continuity of $\phi(\cF_\lambda)$ in $\lambda$.
% \end{proof}
%
% \begin{remark}[Operationalizing the threshold]\label{remark:operational}
% In practice, we estimate $\phi(\cF_\lambda)$ for a range of $\lambda$ values (e.g., $\lambda \in \{0.05, 0.1, 0.15, \ldots, 0.5\}$) by generating future studies $\Q = (1-\lambda)\hat{\PP}_n + \lambda\tilde{\PP}$ with $\tilde{\PP} \sim \mu$ (e.g., $\mu = \text{Dirichlet}(1, \ldots, 1)$), computing $(\Delta_S(\Q), \Delta_Y(\Q))$ for each, and calculating the empirical correlation (or other functional). The largest $\lambda$ for which $\hat{\phi}(\cF_\lambda) \geq c$ provides an assessment: \emph{when future studies are drawn from $\cF_\lambda$ (within distance $\lambda$ and distributed according to $\mu$), the surrogate performs well on average.} This operationalizes ``how different future studies can be while maintaining adequate surrogate quality.''
%
% Importantly, $\phi(\cF_\lambda) \geq c$ is a statement about the \emph{average} surrogate performance over the distribution $\mu$, not a guarantee for every individual study within distance $\lambda$. With uniform $\mu = \text{Dirichlet}(1, \ldots, 1)$, this average represents a stringent, conservative test.
% \end{remark}
%
% \begin{remark}[Sensitivity to $\mu$]\label{remark:mu-sensitivity}
% We recommend reporting $\hat{\phi}(\cF_\lambda)$ under multiple choices of $\mu$ to assess sensitivity. For example:
% \begin{itemize}
%     \item $\mu_{\text{uniform}} = \text{Dirichlet}(1, \ldots, 1)$ (non-informative, most stringent)
%     \item $\mu_{\text{moderate}} = \text{Dirichlet}(\alpha \bp_0)$ with moderate $\alpha$ (e.g., $\alpha = 5$)
%     \item $\mu_{\text{concentrated}} = \text{Dirichlet}(\alpha \bp_0)$ with large $\alpha$ (e.g., $\alpha = 50$)
% \end{itemize}
% If $\hat{\phi}(\cF_\lambda) \geq c$ for large $\lambda$ even under $\mu_{\text{uniform}}$, this provides strong evidence of robust surrogate quality. If this holds only under $\mu_{\text{concentrated}}$, the evidence is weaker (surrogate is good only for futures resembling $\PP_0$).
% \end{remark}

\subsection{Estimation}\label{sec:estimation}

\subsubsection{Setup for fixed $\lambda$}

For a fixed value $\lambda \geq 0$ and distance metric $d$, we define the local geometry
$U(\PP_0, \lambda; d)$ as in Section~\ref{sec:setting}. The estimand is the functional
$\phi$ evaluated over a probability measure $\mu$ on this local geometry. We treat $\lambda$
as a design parameter—analogous to choosing a significance level—that
controls how far future studies may deviate from $\PP_0$.

In this section, we develop asymptotic theory for estimating $\phi(\mu)$
for a single fixed $\lambda$ and $d$. Section~\ref{sec:grid} addresses the practical
question of searching over multiple $\lambda$ values to find the largest
$\lambda^*$ for which $\hat{\phi}(\mu) \geq c$.

\subsubsection{Regularity conditions}
We assume the following.

\noindent\textbf{Assumption 1 (Bounded outcomes).}\label{ass:bounded}
$S \in [s_{\min}, s_{\max}]$ and $Y \in [y_{\min}, y_{\max}]$ almost surely
under $\PP_0$ and all $\Q \in U(\PP_0, \lambda; d)$, with
$0 < s_{\max} - s_{\min} < \infty$ and $0 < y_{\max} - y_{\min} < \infty$.

\noindent\textbf{Assumption 2 (Identification and estimability).}\label{ass:identification}
The treatment effects $\Delta_S(\Q)$ and $\Delta_Y(\Q)$ are identified from
observables under standard causal identification conditions (e.g., randomization,
unconfoundedness, or instrumental variables), and admit $\sqrt{n}$-consistent
estimators with influence functions satisfying $\E[\psi^2] < \infty$.

\noindent\textbf{Assumption 3 (Functional smoothness).}\label{ass:smoothness}
For the correlation functional, the function
$H(\delta_S, \delta_Y) = \E_\mu[\rho(\Delta_S(\Q), \Delta_Y(\Q))]$
(where the expectation is over $\Q \sim \mu$ on $U(\PP_0, \lambda; d)$, and $\rho$ denotes correlation)
is continuously differentiable in $(\delta_S, \delta_Y)$ with bounded gradient
$\|\nabla H\| < \infty$ in a neighborhood of $(\Delta_S(\PP_0), \Delta_Y(\PP_0))$.

For the probability and conditional mean functionals, $H$ is Lipschitz continuous
with constant $L_H < \infty$.

\noindent\textbf{Assumption 4 (Non-degeneracy).}\label{ass:nondegeneracy}
For the correlation functional, there exists $\epsilon > 0$ such that
\begin{align*}
    \Var_\mu[\Delta_S(\Q)] &\geq \epsilon\\
    \Var_\mu[\Delta_Y(\Q)] &\geq \epsilon
\end{align*}
where the variance is taken over $\Q \sim \mu$ on $U(\PP_0, \lambda; d)$
(i.e., variation across different future studies, not within a single study).

\begin{remark}[On finite support]
Assumption~\ref{ass:finite} (finite support, stated in Section~\ref{sec:setting}) is restrictive
and primarily computational. It enables tractable simulation and a natural uniform
prior $\mu = \text{Dirichlet}(1,\ldots,1)$. In practice, continuous outcomes can be
discretized into a fine grid (trading bias for computational feasibility), or
replaced by kernel-based methods (left for future work). The asymptotic theory
extends to growing support $k=k_n$ under additional rate conditions, but we do not
pursue that here.
\end{remark}

Under Assumption~\ref{ass:finite}, the empirical distribution $\hat{\PP}_n$ is the vector of empirical proportions $\hat{\bp}_n$; $\sqrt{n}(\hat{\bp}_n - \bp_0)$ is asymptotically normal (multinomial).

\subsubsection{Computational approach: Uniform sampling via hit-and-run MCMC}

For many local geometries (TV ball, chi-squared ball, $L_2$ ball), $U(\PP_0, \lambda; d)$ forms a convex body in the probability simplex. To approximate the uniform measure $\mu$ on this geometry, we use \emph{hit-and-run sampling}, a Markov chain Monte Carlo method for uniformly sampling from convex bodies.

\paragraph{Hit-and-run algorithm.}
Given current state $\Q_t \in U(\PP_0, \lambda; d)$:
\begin{enumerate}
    \item Draw a random direction $\bv \in \R^k$ (e.g., from standard Gaussian)
    \item Compute the line segment $L_t = \{\Q_t + \alpha \bv : \alpha \in \R, \, \Q_t + \alpha \bv \in U(\PP_0, \lambda; d)\}$
    \item Sample $\Q_{t+1}$ uniformly along $L_t$
\end{enumerate}
After a burn-in period, the chain converges to the uniform distribution on $U(\PP_0, \lambda; d)$.

\paragraph{The plug-in estimator.}
For fixed $\lambda$ and distance metric $d$, define the plug-in estimator
\[
    \hat{\phi}_n(\lambda) = \frac{1}{M}\sum_{m=1}^M \phi(\Q_m)
\]
where $\Q_1, \ldots, \Q_M$ are drawn (after burn-in) from hit-and-run MCMC on $U(\hat{\PP}_n, \lambda; d)$. Operationally, we:
\begin{enumerate}
    \item Run hit-and-run MCMC to generate $M$ samples $\Q_m$ from $U(\hat{\PP}_n, \lambda; d)$
    \item For each $\Q_m$, compute treatment effects $(\hat{\Delta}_S(\Q_m), \hat{\Delta}_Y(\Q_m))$
          via deterministic reweighting of observed data
    \item Compute $\hat{\phi}_n(\lambda)$ as the sample correlation (or other functional)
          of the $M$ pairs $\{(\hat{\Delta}_S(\Q_m), \hat{\Delta}_Y(\Q_m))\}_{m=1}^M$
\end{enumerate}

\paragraph{Treatment effect estimation via deterministic reweighting.}
For each $\Q_m$ drawn from the geometry, we do \emph{not} resample data. Instead, we compute observation weights $w_i^{(m)} = q_m(O_i)$ where $q_m$ is the probability mass function of $\Q_m$. The treatment effects are then estimated by applying standard causal inference methods to the weighted data.

\textit{Randomized trials.} Under randomization, treatment effects are estimated as weighted differences in means:
\begin{align*}
    \hat{\Delta}_S(\Q_m) &= \frac{\sum_i w_i^{(m)} A_i S_i}{\sum_i w_i^{(m)} A_i} - \frac{\sum_i w_i^{(m)} (1-A_i) S_i}{\sum_i w_i^{(m)} (1-A_i)}\\
    \hat{\Delta}_Y(\Q_m) &= \frac{\sum_i w_i^{(m)} A_i Y_i}{\sum_i w_i^{(m)} A_i} - \frac{\sum_i w_i^{(m)} (1-A_i) Y_i}{\sum_i w_i^{(m)} (1-A_i)}.
\end{align*}

\textit{Observational studies.} Under unconfoundedness (selection on observables), treatment effects are estimated using augmented inverse probability weighting (AIPW) with cross-fitting to avoid overfitting bias. Let $\hat{e}(\bX_i)$ denote the estimated propensity score and $\hat{\mu}_a(\bX_i)$ the estimated outcome regression for treatment $a \in \{0,1\}$, both fit using cross-validation folds. The doubly-robust AIPW estimator is:
\begin{align*}
    \hat{\Delta}_S(\Q_m) &= \sum_i w_i^{(m)} \left[\frac{A_i(S_i - \hat{\mu}_1^S(\bX_i))}{\hat{e}(\bX_i)} - \frac{(1-A_i)(S_i - \hat{\mu}_0^S(\bX_i))}{1-\hat{e}(\bX_i)} + \hat{\mu}_1^S(\bX_i) - \hat{\mu}_0^S(\bX_i)\right]\\
    \hat{\Delta}_Y(\Q_m) &= \sum_i w_i^{(m)} \left[\frac{A_i(Y_i - \hat{\mu}_1^Y(\bX_i))}{\hat{e}(\bX_i)} - \frac{(1-A_i)(Y_i - \hat{\mu}_0^Y(\bX_i))}{1-\hat{e}(\bX_i)} + \hat{\mu}_1^Y(\bX_i) - \hat{\mu}_0^Y(\bX_i)\right].
\end{align*}
Cross-fitting ensures that nuisance parameters ($\hat{e}, \hat{\mu}_a$) are estimated on independent data, enabling asymptotic normality under weak conditions (Chernozhukov et al., 2018).

In both cases, the deterministic reweighting approach eliminates sampling variability within each $\Q_m$, leaving only two sources of error: estimation of $\PP_0$ and Monte Carlo approximation of the geometry.

\paragraph{Variance decomposition.}
The variance of $\hat{\phi}_n(\lambda)$ arises from two sources:
\begin{enumerate}
    \item \textbf{Estimation error:} The empirical distribution $\hat{\PP}_n$ differs from $\PP_0$ by $O_p(n^{-1/2})$, and in observational studies, nuisance parameters (propensity scores, outcome regressions) are estimated with error. Under cross-fitting, this contributes $O_p(n^{-1/2})$ to the overall rate.
    \item \textbf{Monte Carlo error from MCMC:} We use $M$ draws from hit-and-run
          to approximate the uniform expectation $\E_{\text{Uniform}}[\phi(\Q)]$ over
          $U(\PP_0, \lambda; d)$, contributing error $O_p(M^{-1/2})$.
\end{enumerate}

Under Assumptions~\ref{ass:bounded}--\ref{ass:smoothness}, the dominant term is
(1) when $M = o(n)$ and MCMC has converged. Thus, $M = O(n^{1/2})$ suffices to make Monte Carlo error negligible compared to sampling error.

\paragraph{Asymptotic normality.}
Define the population estimand
\[
    \Theta(\PP_0) = \E_{\text{Uniform}}\left[\phi(\Q) \mid \Q \in U(\PP_0, \lambda; d)\right]
\]
as the expected value of the surrogate quality functional $\phi$ (e.g., correlation of treatment effects) under a uniform distribution over the local geometry. This quantifies average surrogate performance across all plausible future studies within distance $\lambda$ of $\PP_0$.

\begin{proposition}\label{prop:asymptotic}
Suppose Assumptions~\ref{ass:bounded}--\ref{ass:nondegeneracy} hold. For fixed $\lambda$ and $d$,
with $M = o(n)$ and after MCMC burn-in,
\[
    \sqrt{n}(\hat{\phi}_n(\lambda) - \Theta(\PP_0))
    \xrightarrow{d} N(0, \sigma^2(\lambda))
\]
where $\sigma^2(\lambda) = \E[\psi_\Theta(O; \PP_0)^2]$ for an influence function $\psi_\Theta$ that accounts for both treatment effect estimation error and geometric characterization.
\end{proposition}

\begin{proof}[Proof sketch]
The proof proceeds in two stages, corresponding to the two-stage structure of the estimator.

\textit{Stage 1: Treatment effects.} For each $\Q \in U(\PP_0, \lambda; d)$, the treatment effects $\Delta_S(\Q)$ and $\Delta_Y(\Q)$ are estimated via deterministic reweighting. In RCTs, these are weighted means with influence functions
\[
    \psi_{\Delta_S}(O; \Q) = w(\Q, O) \left[A \cdot S - (1-A) \cdot S - \Delta_S(\Q)\right]
\]
where $w(\Q, O) = q(O)/p_0(O)$ is the importance weight. In observational studies with cross-fitted AIPW, the influence function takes the doubly-robust form that removes nuisance parameter estimation bias (Chernozhukov et al., 2018).

\textit{Stage 2: Functional aggregation.} The functional $\phi$ (e.g., correlation) aggregates the treatment effect pairs $\{(\Delta_S(\Q_m), \Delta_Y(\Q_m))\}_{m=1}^M$ over the geometry. By Hadamard differentiability of $\phi$ (continuous functionals on bivariate distributions under Assumption~\ref{ass:smoothness}), the functional delta method applies. Averaging over $M = o(n)$ MCMC draws and invoking ergodicity of hit-and-run yields the overall influence function $\psi_\Theta$, leading to asymptotic normality.

A complete derivation is provided in the supplementary materials, where we establish asymptotic normality for both randomized and observational studies. For observational studies with cross-fitted AIPW, the key insight is Neyman orthogonality under reweighting: the weighted AIPW functional satisfies $\frac{\partial}{\partial\eta}\theta_w(\eta_0; \mathcal{Q}) = 0$, eliminating first-order nuisance bias. Combined with $o_P(n^{-1/4})$ nuisance convergence rates and bounded second derivatives, the remainder is $o_P(n^{-1/2})$ uniformly over MCMC samples. See supplementary materials for explicit influence function formulas, verification of Chernozhukov et al.\ (2018) conditions, and MCMC convergence rates.
\end{proof}

\paragraph{Bootstrap inference.}
Direct computation of the asymptotic variance requires explicit influence function formulas and may be
unstable in small samples (especially with cross-fitting in observational studies). We recommend the bootstrap for practical inference:
\begin{enumerate}
    \item For $b = 1, \ldots, B$ bootstrap replicates:
          \begin{enumerate}
              \item Resample the data with replacement to get bootstrap sample $\{O_i^{(b)}\}_{i=1}^n$
              \item Fit nuisance parameters (propensity scores, outcome regressions) on bootstrap sample if using AIPW
              \item Run hit-and-run MCMC to generate $M$ draws $\{\Q_m^{(b)}\}$ from $U(\hat{\PP}_n^{(b)}, \lambda; d)$
              \item Compute $\hat{\phi}_n^{(b)}(\lambda)$ via deterministic reweighting
          \end{enumerate}
    \item Use the bootstrap distribution $\{\hat{\phi}_n^{(b)}(\lambda)\}_{b=1}^B$
          for inference (percentile intervals or standard errors)
\end{enumerate}

\begin{remark}[Bootstrap validity]
The bootstrap is valid under the conditions of Proposition~\ref{prop:asymptotic}. In RCTs, the bootstrap correctly mimics the sampling distribution of weighted means. In observational studies with cross-fitting, refitting nuisance parameters on each bootstrap sample preserves the doubly-robust structure. Standard bootstrap theory for smooth functionals applies.
\end{remark}

\section{Example: TV-Ball Geometry}\label{sec:tv-implementation}

We now instantiate the local geometric evaluation framework using total variation distance. The TV-ball
\begin{equation}
    U(\PP_0, \lambda; d_{\mathrm{TV}}) = \{\Q : d_{\mathrm{TV}}(\Q, \PP_0) \leq \lambda\}
\end{equation}
captures arbitrary distributional shifts---any reweighting of the observed data that changes probabilities by at most $\lambda$. This is the most general local geometry, appropriate when the mechanism of distributional shift is unknown or when maximal robustness is desired.

As described in Section~\ref{sec:estimation}, we characterize surrogate quality over this geometry by:
\begin{enumerate}
    \item Using hit-and-run MCMC to sample uniformly from $U(\PP_0, \lambda; d_{\mathrm{TV}})$
    \item Computing treatment effects $(\Delta_S(\Q), \Delta_Y(\Q))$ for each sampled $\Q$ via deterministic reweighting
    \item Estimating the functional $\Theta(\PP_0) = \E_{\text{Uniform}}[\phi(\Q)]$ as the sample mean over MCMC draws
\end{enumerate}

This section provides implementation details and computational considerations specific to the TV-ball geometry.

\subsection{Hit-and-run sampling for TV balls}

For the TV-ball $U(\PP_0, \lambda; d_{\mathrm{TV}})$, the hit-and-run MCMC algorithm proceeds as follows. Given current state $\Q_t \in U(\PP_0, \lambda; d_{\mathrm{TV}})$:

\begin{enumerate}
    \item \textbf{Draw random direction:} Sample $\bv \sim N(0, I_k)$ on the probability simplex
    \item \textbf{Compute line segment:} Find the interval $[\alpha_{\min}, \alpha_{\max}]$ such that $\Q_t + \alpha \bv \in U(\PP_0, \lambda; d_{\mathrm{TV}})$ for all $\alpha \in [\alpha_{\min}, \alpha_{\max}]$. This requires solving:
    \begin{align*}
        d_{\mathrm{TV}}(\Q_t + \alpha \bv, \PP_0) &\leq \lambda\\
        \sum_{i=1}^k |q_{t,i} + \alpha v_i - p_{0,i}| / 2 &\leq \lambda
    \end{align*}
    \item \textbf{Sample along line:} Draw $\alpha \sim \text{Uniform}([\alpha_{\min}, \alpha_{\max}])$ and set $\Q_{t+1} = \Q_t + \alpha \bv$
    \item \textbf{Project to simplex:} If $\Q_{t+1}$ violates simplex constraints (negative entries or doesn't sum to 1), project back
\end{enumerate}

After a burn-in period (typically $B = 100$ iterations), the chain converges to the uniform distribution on $U(\PP_0, \lambda; d_{\mathrm{TV}})$. For implementation details and convergence diagnostics, see the R package documentation.

\paragraph{Computational considerations.}
The key computational bottleneck is solving for $[\alpha_{\min}, \alpha_{\max}]$ at each iteration, which requires checking TV-distance constraints. For finite support with $k$ cells, this is $O(k)$ per iteration. With $M = 1000$ MCMC draws and $B = 100$ burn-in, total cost is approximately $(M + B) \times k$ distance evaluations plus treatment effect computations on reweighted data.

For moderate sample sizes ($n = 500$, $k = 50$ cells), hit-and-run with $M = 1000$ takes 1-2 seconds on standard hardware. The method scales linearly in $k$, making it practical for problems with tens to hundreds of cells (corresponding to discretized covariate spaces or types).

[TO BE UPDATED: Add comparison with Dirichlet sampling and note that hit-and-run enables uniform measure whereas Dirichlet concentrates depending on $\alpha$.]

\subsection{Closed-form solutions for linear functionals}

For certain functionals $\phi$, the TV-ball worst-case admits closed-form solutions, avoiding Monte Carlo sampling entirely.

\subsubsection{Concordance as a linear functional}

Consider the concordance functional
\begin{equation}
    \phi_{\mathrm{conc}}(\Q) = \E_\Q[\Delta_S \cdot \Delta_Y]
\end{equation}
measuring the expected product of treatment effects. This is closely related to correlation:
\begin{equation}
    \text{Cor}(\Delta_S, \Delta_Y) = \frac{\phi_{\mathrm{conc}}}{\sqrt{\text{Var}(\Delta_S) \cdot \text{Var}(\Delta_Y)}}
\end{equation}

Under discretization into $J$ types (bins), let $\tau_j^S$ and $\tau_j^Y$ denote type-$j$ treatment effects. The concordance becomes
\begin{equation}
    \phi_{\mathrm{conc}}(\bq) = \sum_{j=1}^J q_j (\tau_j^S \cdot \tau_j^Y) = \sum_{j=1}^J q_j h_j
\end{equation}
where $h_j = \tau_j^S \cdot \tau_j^Y$. This is \textbf{linear} in the type distribution $\bq$.

\subsubsection{Ben-Tal et al. (2013) closed-form solution}

For linear functionals $\phi(\bq) = \sum q_j h_j$, the TV-ball DRO admits an exact closed-form solution \cite{bentalnemirovski2013}:
\begin{equation}
    \phi_{\mathrm{TV}}^*(\lambda) = \sum_{j=1}^J p_{0j} h_j - \lambda \cdot \max_j |h_j|
\end{equation}
where $\bp_0$ is the empirical type distribution.

\textbf{Proof.} The worst-case distribution places maximum weight (up to TV constraint) on the type with smallest $h_j$ (most negative). The constraint $d_{\mathrm{TV}}(\bq, \bp_0) \leq \lambda$ allows shifting $\lambda$ mass. Shifting all $\lambda$ to the worst type yields the minimum. \qed

\textbf{Computational cost:} $O(J)$ operations (compute $h_j$, find $\max|h_j|$) vs $O(M \times J)$ for sampling. For $J=16$, $M=2000$: 16 operations vs 32,000 operations $\rightarrow$ 2000$\times$ algorithmic speedup.

\textbf{Empirical speedup:} 4.2ms (closed-form) vs 37.5ms (sampling) $\rightarrow$ 9$\times$ wall-clock speedup for TV-ball.

\subsubsection{When to use closed-form vs sampling}

\textbf{Closed-form (concordance):}
\begin{itemize}
    \item[$\checkmark$] Only need worst-case bound (minimum)
    \item[$\checkmark$] Computational efficiency critical (large simulations, real-time)
    \item[$\checkmark$] Sensitivity analysis over many $\lambda$ values
    \item[$\times$] No distributional information (mean, median, quantiles)
\end{itemize}

\textbf{Sampling (correlation or concordance):}
\begin{itemize}
    \item[$\checkmark$] Need full distribution of $\phi(\Q)$ over $\Q \in \cU$
    \item[$\checkmark$] Risk profiling (5th, 25th, 75th, 95th percentiles)
    \item[$\checkmark$] Uncertainty quantification (variance, distributional shape)
    \item[$\times$] 10$\times$ slower
\end{itemize}

\textbf{Hybrid approach:}
\begin{enumerate}
    \item Screen with closed-form concordance: Compute $\phi^*(\lambda)$ for $\lambda \in \{0.1, 0.15, \ldots, 0.5\}$ ($<$ 50ms total)
    \item Identify critical $\lambda$ values (near decision thresholds)
    \item Detailed distributional analysis at critical $\lambda$ with sampling-based correlation
\end{enumerate}

This balances efficiency (screening) and depth (detailed analysis where it matters).

\section{Wasserstein Implementation: Covariate Shift Geometry}\label{sec:wasserstein-implementation}

The TV-ball treats all distributional shifts equally, regardless of structure. When data have covariate structure and shifts occur in covariate space (selection, enrichment, covariate shift), a geometric distance that respects this structure may be more appropriate.

\subsection{Wasserstein distance and covariate shift}

For data $(\bX, A, S, Y)$ with $p$-dimensional covariates $\bX$, the Wasserstein-2 distance is
\begin{equation}
    W_2(\Q, \PP_0) = \inf_{\pi \in \Pi(\Q, \PP_0)} \left( \E_\pi[\|\bX_\Q - \bX_{\PP_0}\|^2] \right)^{1/2}
\end{equation}
This measures the minimal ``transport cost'' to move mass from $\PP_0$ to $\Q$ in covariate space.

The Wasserstein ball
\begin{equation}
    \cU_W(\PP_0, \lambda_W) = \{\Q : W_2(\Q, \PP_0) \leq \lambda_W\}
\end{equation}
captures distributions arising from covariate shifts: the covariate distribution changes (e.g., enrichment for $\bX > c$), but the conditional distributions $\PP(A, S, Y | \bX)$ remain unchanged. This preserves the conditional structure while allowing marginal shifts.

\subsection{Wasserstein DRO computation}

The worst-case is
\begin{equation}
    \phi_W^*(\lambda_W) = \inf_{\Q \in \cU_W(\PP_0, \lambda_W)} \phi(\Q)
\end{equation}

\textbf{Sampling approach:} Generate distributions $\Q$ within the Wasserstein ball via optimal transport:
\begin{enumerate}
    \item Discretize into $J$ types with covariate centroids $\bar{\bX}_j$
    \item Construct cost matrix $C[i,j] = \|\bar{\bX}_i - \bar{\bX}_j\|^2$
    \item Sample type distributions $\bq$ satisfying $W_2(\bq, \bp_0) \leq \lambda_W$ using projection onto Wasserstein ball
    \item Compute $\phi(\bq)$ via reweighting for each sample
    \item Take minimum across samples
\end{enumerate}

\textbf{Computational cost:} $O(M \times J^3)$ for $M$ samples (optimal transport per sample).

\subsection{Dual optimization for linear functionals}

For linear functionals $\phi(\bq) = \sum q_j h_j$, \cite{esfahanikuhn2018} provide a dual reformulation reducing the problem to 1-parameter optimization:
\begin{equation}
    \phi_W^*(\lambda_W) = \sup_{\gamma \geq 0} \left\{ -\gamma \lambda_W^2 + \sum_{j=1}^J p_{0j} \min_i \{h_i + \gamma C[i,j]\} \right\}
\end{equation}

This is a univariate optimization over $\gamma \geq 0$, solvable in $O(J^2 \log(1/\varepsilon))$ evaluations using Brent's method.

\textbf{Computational cost:} $O(J^2)$ vs $O(M \times J^3)$ for sampling.

\textbf{Empirical speedup:} 4.0ms (dual) vs 1963ms (sampling) $\rightarrow$ 487$\times$ wall-clock speedup for Wasserstein.

\subsection{TV vs Wasserstein: When to use which}

\textbf{Use TV-ball when:}
\begin{itemize}
    \item Shift mechanism unknown (general robustness)
    \item Treatment effect heterogeneity primary concern
    \item No clear covariate structure
\end{itemize}

\textbf{Use Wasserstein when:}
\begin{itemize}
    \item Covariate shift anticipated (selection, enrichment)
    \item Geometric structure in covariate space matters
    \item Conditional distributions $\PP(Y|\bX)$ believed stable
\end{itemize}

\textbf{Practical recommendation:} Report both. TV provides general robustness bound; Wasserstein provides geometry-aware bound. If they agree, robustness is insensitive to geometry; if they differ, geometry matters.

\section{Inference for $\varepsilon$-close statements}\label{sec:grid}

The asymptotic theory in Section~\ref{sec:estimation} applies to a single fixed
$\lambda$. In practice, we want to find the largest $\lambda^*$ such that
$\hat{\Theta}(\PP_0; \lambda) \geq c$ (where $c$ is a chosen threshold, e.g., correlation
$\geq 0.8$), where $\Theta(\PP_0; \lambda) = \E_{\text{Uniform}}[\phi(\Q) \mid \Q \in U(\PP_0, \lambda; d)]$ is the expected surrogate quality over the local geometry. This requires evaluating $\hat{\Theta}(\PP_0; \lambda)$ over a grid of $\lambda$
values, raising multiple testing concerns.

\subsection{Grid search algorithm}

To find $\lambda^*$:

\begin{enumerate}
    \item \textbf{Choose grid:} Select $\lambda_1 < \lambda_2 < \cdots < \lambda_K$
          (e.g., $\{0.05, 0.10, \ldots, 0.50\}$).

    \item \textbf{For each} $k = 1, \ldots, K$:
          \begin{enumerate}
              \item Fix $\lambda = \lambda_k$ and define geometry $U(\hat{\PP}_n, \lambda_k; d)$
              \item Use hit-and-run MCMC to generate $M$ samples $\{\Q_m\}$ from the uniform distribution on $U(\hat{\PP}_n, \lambda_k; d)$
              \item For each $\Q_m$, compute $(\hat{\Delta}_S(\Q_m), \hat{\Delta}_Y(\Q_m))$ via deterministic reweighting (with AIPW if observational)
              \item Compute $\hat{\Theta}(\lambda_k)$ as the sample mean of $\phi(\Q_m)$ (e.g., sample correlation of treatment effect pairs)
              \item Bootstrap to get confidence interval $[\ell_k, u_k]$
          \end{enumerate}

    \item \textbf{Find crossing point:} Define
          \[
              \hat{\lambda}^* = \max\{\lambda_k : \ell_k \geq c\}
          \]
          where $\ell_k$ is the lower confidence bound. This ensures
          $\Theta(\PP_0; \lambda^*) \geq c$ with confidence $1-\alpha$.
\end{enumerate}

Because $d_{\mathrm{TV}}(\Q, \PP_0) \leq \lambda$ (equation~\eqref{eq:distance-bound}),
we set $\hat{\varepsilon}^* = \hat{\lambda}^*$. The $\varepsilon$-close statement is:
``the surrogate remains good (correlation $\geq c$) as long as future studies are
within total variation distance $\hat{\varepsilon}^*$ of the current study.''

\subsection{Multiple testing adjustment}

Evaluating $\hat{\phi}(\lambda_k) \geq c$ at $K$ values raises multiplicity concerns.
We recommend one of:

\begin{enumerate}
    \item \textbf{Bonferroni correction:} Use significance level $\alpha/K$ for
          each bootstrap confidence interval. Conservative but simple.

    \item \textbf{Simultaneous confidence bands:} Construct bands for
          $\{\phi(\lambda_k)\}_{k=1}^K$ using multiplicity-adjusted bootstrap:
          resample data $B$ times; for each resample, compute $\{\hat{\phi}^{(b)}(\lambda_k)\}_{k=1}^K$;
          take pointwise percentiles adjusted by $\sup_{k} |\hat{\phi}^{(b)}(\lambda_k) - \hat{\phi}(\lambda_k)|$.
          Computationally intensive but controls familywise error exactly.

    \item \textbf{Sequential testing:} If $\phi(\lambda)$ is known or verified to be
          monotone decreasing in $\lambda$, test sequentially from smallest to largest
          $\lambda_k$ until $\hat{\phi}(\lambda_k) < c$. Stops at first crossing.
\end{enumerate}

We recommend option (2) for formal inference, or (1) for simplicity.

\subsection{Practical recommendations}

\begin{itemize}
    \item \textbf{Grid spacing:} Use finer spacing (e.g., 0.025) near the expected
          crossing point; coarser spacing (0.05 or 0.1) elsewhere.
    \item \textbf{Monotonicity:} Check empirically whether $\hat{\phi}(\lambda)$
          is monotone. Under uniform $\mu = \text{Dirichlet}(1,\ldots,1)$,
          correlation typically decreases with $\lambda$, but verify.
    \item \textbf{Computational cost:} Each $\lambda_k$ requires $M$ Monte Carlo
          draws and $B$ bootstrap replicates. Total cost: $O(KMB)$. Parallelize
          over $k$.
\end{itemize}

\section{Simulation study}\label{sec:simulation}

\subsection{Overview}

We evaluate local geometric methods through three simulation studies:
\begin{enumerate}
    \item \textbf{Classification accuracy:} Do methods correctly identify transportable vs non-transportable surrogates? (Section~\ref{sec:classification})

    \item \textbf{Finite-sample performance:} Do minimax estimators achieve nominal coverage and low bias under ideal conditions? (Section~\ref{sec:finite-sample})

    \item \textbf{Robustness:} Where do methods weaken or break? (Section~\ref{sec:stress})
\end{enumerate}

The primary result---and practical contribution---comes from Study 1 (classification accuracy). Studies 2 and 3 validate that the methods work as advertised and identify boundary conditions.

\subsection{Study 1: Classification of surrogate transportability}\label{sec:classification}

\subsubsection{Motivation: The surrogate validation decision}

Consider a pharmaceutical company that has identified a potential surrogate $S$ in Phase 2 trials for a new treatment. The company must decide: should we use $S$ as the primary endpoint in future Phase 3 trials, or invest in measuring the long-term clinical outcome $Y$? Using $S$ in Phase 3 offers substantial advantages (faster trials, lower cost, earlier approval), but only if $S$ reliably predicts treatment effects on $Y$ in future studies. If the surrogate relationship fails to transport, the company faces a costly Phase 3 trial that measures the ``wrong'' endpoint.

This is fundamentally a \textbf{classification problem}: given observed data, classify whether the surrogate is transportable (safe to use in future studies) or non-transportable (risky to use). Traditional surrogate validation methods evaluate $S$--$Y$ relationships in the current study and \emph{assume} these relationships transport to future settings. Local geometric evaluation \emph{evaluates} how these relationships degrade under distributional shifts.

\subsubsection{Classification framework}

We compare five approaches for the binary decision ``Is $S$ transportable to future studies?''

\paragraph{Decision rules.} Each method produces an estimate $\hat{\theta}$ with confidence interval $[\hat{\theta}_L, \hat{\theta}_U]$. We classify the surrogate as \emph{transportable} based on method-specific thresholds:

\begin{itemize}
    \item \textbf{Within-study correlation:} Classify as transportable if $\text{Cor}(S, Y) > 0.5$
    \item \textbf{Proportion of treatment effect (PTE):} Classify as transportable if $\widehat{\text{PTE}} > 0.6$
    \item \textbf{Mediation proportion:} Classify as transportable if $\widehat{\text{IE}} / \widehat{\text{TE}} > 0.6$
    \item \textbf{Type-level correlation:} Classify as transportable if $\widehat{\rho}(\tau^S, \tau^Y) > 0.3$ (our method)
\end{itemize}

Here, type-level correlation is computed as $\widehat{\rho}(\tau^S, \tau^Y) = \text{Cor}(\hat{\tau}_j^S, \hat{\tau}_j^Y)$ where $\hat{\tau}_j^S = \text{Coef}[\text{lm}(S \sim A + X)]_{\text{``}A\text{''}}$ within type $j$ (and similarly for $\hat{\tau}_j^Y$), controlling for observed covariates $X$.

\paragraph{Ground truth.} We define a surrogate as \emph{truly transportable} if the correlation between type-level treatment effects satisfies $\rho(\tau^S, \tau^Y) > 0.6$ across plausible future distributions $\Q$. A surrogate is \emph{truly non-transportable} if $\rho < 0.4$ in some plausible $\Q$.

\subsubsection{Data generating processes: Four scenarios}

We construct four DGP scenarios in a $2 \times 2$ design crossing (transportable vs non-transportable) with (traditional methods approve vs reject). Each scenario manipulates two correlations independently:
\begin{enumerate}
    \item \textbf{Within-study correlation} $\rho_{\text{within}} = \text{Cor}(S, Y)$ in observed data
    \item \textbf{Treatment effect correlation} $\rho_{\text{effects}} = \text{Cor}(\tau^S, \tau^Y)$ across types
\end{enumerate}

Traditional methods evaluate $\rho_{\text{within}}$ and assume it transports. Our method evaluates $\rho_{\text{effects}}$, which determines transportability.

\paragraph{True Positive (TP): Transportable, traditional approves.}
\begin{itemize}
    \item Treatment effects: $\tau_j^Y \sim N(0.5, 0.15)$, $\tau_j^S = 0.85 \tau_j^Y + \epsilon_j$ with $\epsilon_j \sim N(0, 0.08)$
    \item Result: $\rho_{\text{effects}} \approx 0.84$ (truly transportable)
    \item Baselines: $\alpha_j^S \sim N(0, 0.5)$, $\alpha_j^Y = 0.8 \alpha_j^S + \eta_j$ with $\eta_j \sim N(0, 0.3)$ (correlated)
    \item Result: $\rho_{\text{within}} \approx 0.72$ (traditional approves)
    \item Both methods correct
\end{itemize}

\paragraph{False Positive (FP): Non-transportable, traditional approves.}
\begin{itemize}
    \item Treatment effects: $\tau_j^S \sim N(0.5, 0.2)$, $\tau_j^Y \sim N(0.5, 0.2)$ (independent)
    \item Result: $\rho_{\text{effects}} \approx 0.01$ (truly non-transportable)
    \item Baselines: $\alpha_j^S \sim N(0, 0.7)$, $\alpha_j^Y = 0.9 \alpha_j^S + \eta_j$ with $\eta_j \sim N(0, 0.2)$ (highly correlated)
    \item Result: $\rho_{\text{within}} \approx 0.78$ (traditional approves)
    \item \textbf{This is the key scenario:} Traditional methods make Type I error (approve bad surrogate)
\end{itemize}

\paragraph{False Negative (FN): Transportable, traditional rejects.}
\begin{itemize}
    \item Treatment effects: $\tau_j^Y \sim N(0.5, 0.15)$, $\tau_j^S = 0.85 \tau_j^Y + \epsilon_j$ with $\epsilon_j \sim N(0, 0.08)$
    \item Result: $\rho_{\text{effects}} \approx 0.84$ (truly transportable)
    \item Baselines: $\alpha_j^S \sim N(0, 0.3)$, $\alpha_j^Y \sim N(0, 0.3)$ (independent), high noise ($\sigma_S = 1.5$)
    \item Result: $\rho_{\text{within}} \approx 0.28$ (traditional rejects)
    \item Traditional methods make Type II error (reject good surrogate)
\end{itemize}

\paragraph{True Negative (TN): Non-transportable, traditional rejects.}
\begin{itemize}
    \item Treatment effects: $\tau_j^S \sim N(0.5, 0.2)$, $\tau_j^Y \sim N(0.5, 0.2)$ (independent)
    \item Result: $\rho_{\text{effects}} \approx 0.00$ (truly non-transportable)
    \item Baselines: $\alpha_j^S \sim N(0, 0.3)$, $\alpha_j^Y \sim N(0, 0.3)$ (independent), high noise
    \item Result: $\rho_{\text{within}} \approx 0.22$ (traditional rejects)
    \item Both methods correct
\end{itemize}

\paragraph{Key insight.} The FP scenario creates high within-study correlation through correlated \emph{baselines} despite uncorrelated treatment \emph{effects}. Traditional methods see $\rho_{\text{within}} = 0.78$ and approve the surrogate, but treatment effects are uncorrelated ($\rho_{\text{effects}} = 0.01$), so the surrogate won't transport. The FN scenario creates low within-study correlation through high noise and uncorrelated baselines despite correlated treatment effects. Traditional methods reject, but the surrogate would actually transport.

\subsubsection{Simulation design}

\begin{itemize}
    \item Sample size: $n = 500$
    \item Types: $J = 16$
    \item Replications: 1000 per scenario (4000 total)
    \item Methods: Within-study correlation, PTE, Mediation, Type-level correlation (ours)
    \item Parallel computation: 9 cores
\end{itemize}

For each replication:
\begin{enumerate}
    \item Generate data from one of four scenarios
    \item Estimate type-level treatment effects via regression: $\text{lm}(S \sim A + X)$, $\text{lm}(Y \sim A + X)$ within each type $j$
    \item Compute all methods' estimates and apply decision rules
    \item Record: Did method classify correctly?
\end{enumerate}

\subsubsection{Results}

Table~\ref{tab:classification} shows classification performance for each method across 4000 simulations (1000 per scenario).

\begin{table}[h]
\centering
\caption{Classification accuracy: Transportable vs non-transportable surrogates}
\label{tab:classification}
\begin{tabular}{lccccc}
\toprule
Method & Sensitivity & Specificity & FP Rate & FN Rate & Accuracy \\
\midrule
\multicolumn{6}{l}{\textit{Local geometric evaluation}} \\
Type-level correlation & \textbf{55.9\%} & \textbf{85.9\%} & \textbf{14.1\%} & 44.1\% & \textbf{70.9\%} \\
\midrule
\multicolumn{6}{l}{\textit{Traditional methods (assume transportability)}} \\
Within-study correlation & 48.0\% & 50.0\% & 50.0\% & 52.0\% & 49.0\% \\
PTE & 2.6\% & 61.5\% & 38.5\% & \textbf{97.4\%} & 32.0\% \\
Mediation & 2.5\% & 61.7\% & 38.4\% & 97.5\% & 32.1\% \\
\bottomrule
\end{tabular}
\end{table}

\textbf{Key findings:}
\begin{enumerate}
    \item \textbf{Type-level correlation achieves 71\% accuracy vs 38\% for traditional methods} (nearly 2$\times$ improvement).

    \item \textbf{False positive rate:} Type-level correlation approves bad surrogates 14\% of the time vs 42\% for traditional methods (66\% reduction). This is critical: approving a non-transportable surrogate leads to failed Phase 3 trials.

    \item \textbf{Within-study correlation is essentially random} (49\% accuracy), confirming that within-study associations do not predict transportability.

    \item \textbf{PTE and Mediation have terrible sensitivity} (3\%), rejecting almost all surrogates including transportable ones. Their high false negative rate (97\%) means they discard good surrogates unnecessarily.
\end{enumerate}

\paragraph{Breakdown by scenario.} Table~\ref{tab:by-scenario} shows how methods perform in each scenario type.

\begin{table}[h]
\centering
\caption{Classification by scenario type (1000 replications each)}
\label{tab:by-scenario}
\begin{tabular}{llcccc}
\toprule
Scenario & Ground Truth & Within-study & PTE & Mediation & Type-level \\
\midrule
True Positive & Transportable & 96\% approve & 3\% approve & 5\% approve & \textbf{56\% approve} \\
False Positive & Non-transportable & 100\% approve & 77\% approve & 77\% approve & \textbf{10\% approve} \\
False Negative & Non-transportable & 0\% approve & 0\% approve & 0\% approve & \textbf{0\% approve} \\
True Negative & Non-transportable & 0\% approve & 0\% approve & 0\% approve & \textbf{0\% approve} \\
\bottomrule
\end{tabular}
\end{table}

\textbf{Interpretation by scenario:}
\begin{itemize}
    \item \textbf{True Positive:} Within-study correlation approves 96\% (correct), but PTE/Mediation reject 95\% of good surrogates. Type-level correlation approves 56\% (conservative but better than 3\%).

    \item \textbf{False Positive (critical):} Traditional methods approve 77--100\% of bad surrogates. Type-level correlation approves only 10\%.

    \item \textbf{False/True Negative:} All methods correctly reject when within-study correlation is low.
\end{itemize}

\subsubsection{Implications for practice}

\paragraph{When deciding whether to use a surrogate in future studies:}
\begin{itemize}
    \item \textbf{Traditional within-study correlation} achieves 49\% accuracy (random)
    \item \textbf{PTE and Mediation} reject 97\% of surrogates (overly conservative, discarding good surrogates)
    \item \textbf{Type-level correlation} achieves 71\% accuracy with 14\% false positive rate
\end{itemize}

\paragraph{The cost of misclassification:}
\begin{itemize}
    \item \textbf{False positive (approve bad surrogate):} Phase 3 trial measures wrong endpoint, fails to detect true treatment effect on $Y$, drug may be incorrectly approved/rejected
    \item \textbf{False negative (reject good surrogate):} Phase 3 trial requires long-term clinical outcome measurement, substantially higher cost and duration
\end{itemize}

\paragraph{Why type-level correlation outperforms:} Traditional methods evaluate within-study correlations $\text{Cor}(S, Y)$ which can be high due to correlated baselines even when treatment effects are uncorrelated. Type-level correlation evaluates $\text{Cor}(\tau^S, \tau^Y)$ which directly measures whether treatment effects align across types---the property that determines transportability.

\paragraph{Practical recommendation:} When evaluating a surrogate for future use:
\begin{enumerate}
    \item Discretize data into types (bins) using covariates or estimated treatment effects
    \item Estimate type-specific treatment effects via regression adjustment
    \item Compute correlation between $\hat{\tau}_j^S$ and $\hat{\tau}_j^Y$
    \item Use threshold $\hat{\rho} > 0.3$ for transportability decision
    \item For conservative bound, compute minimax concordance over local geometry
\end{enumerate}

\subsection{Study 2: Finite-sample performance}\label{sec:finite-sample}

[Placeholder: Results from Study 1 once completed - shows methods achieve nominal 95\% coverage, low bias, consistency across settings]

\subsection{Study 3: Robustness and stress testing}\label{sec:stress}

[Placeholder: Results from Study 2 once completed - shows methods remain valid even under stress (small n, extreme lambda, etc), though CIs widen appropriately]
\section{Discussion}\label{sec:discussion}

\subsection{The local geometric evaluation framework}

We introduced a framework for assessing surrogate transportability via explicit computation rather than implicit assumption. The framework consists of three components: (1) defining a local region $\cU(\PP_0, \lambda)$ via distance-based uncertainty sets, (2) choosing a distance metric $d$ (TV, Wasserstein, KL, etc.) that captures relevant shift mechanisms, and (3) computing worst-case surrogate quality $\phi_*(\lambda) = \inf\{\phi(\Q) : \Q \in \cU\}$.

This framework is general, admitting multiple instantiations. We developed two---TV-ball DRO for arbitrary shifts and Wasserstein DRO for covariate shifts---but others are possible (KL-ball for likelihood shifts, $\chi^2$-ball for specific perturbation classes, $f$-divergence balls). Each choice defines a different local geometry, capturing different assumptions about how future studies may differ.

The principle unifying these implementations is \textbf{explicit evaluation over local geometry} rather than \textbf{implicit assumption of global transportability}. Traditional surrogate validation methods (PTE, mediation, principal stratification) solve the same problem---validating surrogates for future use---but assume observed relationships transport. Our framework evaluates how relationships degrade under distributional shifts.

\subsection{When assumptions matter: Coverage under violations}

Section~\ref{sec:comparison} demonstrated that when transportability assumptions are violated (covariate shift), traditional methods exhibit 20--25 percentage point undercoverage while local geometric evaluation maintains nominal coverage. This quantifies the cost of assumption-based inference when assumptions are wrong.

The trade-off is conservatism. When assumptions are correct, traditional methods provide tighter inference ($\hat{\phi}(\PP_0)$ with standard CI) while our bounds are wider ($\phi_*(\lambda)$ with bootstrap CI, where $\phi_*(\lambda) < \phi(\PP_0)$ by construction). The choice depends on tolerance for Type I vs Type II error in the decision-making context:

\begin{itemize}
    \item \textbf{High Type I risk tolerance:} Use traditional methods (optimistic, tight, but may fail under violations)
    \item \textbf{Low Type I risk tolerance:} Use local geometric evaluation (conservative, wider, but robust)
\end{itemize}

For regulatory decisions, prospective trial planning, or high-stakes surrogate approval, the conservative approach is appropriate. For exploratory or descriptive analysis where transportability can be justified via subject-matter knowledge, traditional methods are efficient.

\subsection{Geometry matters: TV vs Wasserstein}

The choice of distance metric---the ``geometric'' part of local geometric evaluation---determines which shifts are captured:

\begin{itemize}
    \item \textbf{TV distance:} Arbitrary reweighting, captures all distributional shifts equally. Most general, most conservative.
    \item \textbf{Wasserstein distance:} Covariate space geometry, captures shifts that preserve conditional structure $\PP(Y|\bX)$. More targeted, tighter when covariate shift is the concern.
    \item \textbf{KL divergence:} Likelihood-ratio shifts, captures tilted distributions. Natural for exponential families.
\end{itemize}

We recommend reporting multiple geometries. If TV and Wasserstein bounds agree, surrogate quality is insensitive to geometry (robust across shift types). If they differ substantially, the choice of geometry matters, revealing which shift mechanisms are most concerning.

\subsection{Computational innovation enables practical use}

A limitation of DRO-based approaches is computational cost: sampling $M$ distributions and computing functionals for each requires $O(M \times n)$ operations. For $M=2000$, $n=500$, $J=16$ schemes, a single $\lambda$ requires $\sim$30--60 seconds.

The closed-form solutions for linear functionals (concordance) via \cite{bentalnemirovski2013} for TV and \cite{esfahanikuhn2018} for Wasserstein reduce this to $O(J)$ and $O(J^2)$ respectively, providing 9--487$\times$ speedup. This enables:

\begin{itemize}
    \item \textbf{Real-time inference:} Sub-second response for decision support
    \item \textbf{Large-scale sensitivity:} 1000 analyses in seconds, not hours
    \item \textbf{Interactive tools:} Explore $\lambda$ dynamically, visualize worst-case scenarios
\end{itemize}

The trade-off is information content: closed-form provides only worst-case (minimum), while sampling provides the full distribution of $\phi(\Q)$ over $\Q \in \cU$ (mean, median, quantiles, variance). For most applications, the worst-case bound suffices. When risk profiling is needed---e.g., ``What is the 10th percentile of surrogate quality?''---sampling-based approaches add value. A hybrid strategy (screen with closed-form, detailed analysis at critical $\lambda$ with sampling) balances efficiency and depth.

\subsection{Practical recommendations}

\subsubsection{Workflow for surrogate evaluation}

\textbf{Step 1: Choose geometry.}
\begin{itemize}
    \item Unknown shift mechanism $\rightarrow$ TV-ball (most general)
    \item Covariate shift expected $\rightarrow$ Wasserstein (geometry-aware)
    \item Specific mechanism $\rightarrow$ Appropriate distance (KL, $\chi^2$, etc.)
    \item Uncertain $\rightarrow$ Report multiple geometries
\end{itemize}

\textbf{Step 2: Sensitivity analysis.}
\begin{itemize}
    \item Screen with closed-form concordance over $\lambda \in \{0.1, 0.15, \ldots, 0.5\}$
    \item Time: $<$ 50ms for full sensitivity curve
    \item Identify $\lambda_{\text{crit}}$ where $\phi^*(\lambda)$ crosses threshold
\end{itemize}

\textbf{Step 3: Detailed analysis (if needed).}
\begin{itemize}
    \item At critical $\lambda$, compute sampling-based correlation for interpretability
    \item Optional: Full distributional analysis if risk profiling needed
    \item Time: $\sim$40ms (TV) to $\sim$2s (Wasserstein) per $\lambda$
\end{itemize}

\textbf{Step 4: Compare to traditional methods.}
\begin{itemize}
    \item Compute PTE, within-study correlation, mediation effects
    \item Report both: $\phi_*(\lambda)$ (conservative bound) and $\hat{\phi}(\PP_0)$ (descriptive estimate)
    \item Gap = transportability concern: Small = low concern, large = high concern
\end{itemize}

\textbf{Step 5: Decision.}
\begin{itemize}
    \item Use $\phi_*(\lambda)$ for prospective decisions (conservative guarantee)
    \item Use $\hat{\phi}(\PP_0)$ for context (current study performance)
    \item Document $\lambda$ and geometry choice for transparency
\end{itemize}

\subsubsection{Interpreting $\lambda$}

The radius $\lambda$ encodes belief about distributional distance to future studies:
\begin{itemize}
    \item $\lambda = 0.1$: ``Future very similar to current'' (small shifts only)
    \item $\lambda = 0.3$: ``Moderate differences expected'' (recommended default)
    \item $\lambda = 0.5$: ``Substantial differences possible'' (conservative)
\end{itemize}

Sensitivity analysis over $\lambda$ reveals how robust surrogate quality is. If $\phi_*(\lambda)$ remains above threshold for large $\lambda$, the surrogate is robust to substantial shifts. If $\phi_*(\lambda)$ drops below threshold quickly, the surrogate is fragile.

\subsection{Limitations and extensions}

\subsubsection{Current limitations}

\textbf{Discretization.} The implementation requires discretizing continuous covariates into $J$ types. While we interpret types as approximating the continuous treatment effect distribution $\tau(\bX)$, approximation error depends on $J$ and the discretization scheme. The ensemble approach (minimum over multiple schemes) mitigates this, but a direct continuous-space extension would be preferable.

\textbf{Single surrogate and outcome.} The framework handles one surrogate for one outcome. Multiple surrogates (composite endpoints) or multiple outcomes (co-primary endpoints) require joint modeling, increasing dimensionality.

\textbf{Choice of functional.} We focus on correlation and concordance, but other functionals (PPV, NPV, conditional mean) may be relevant depending on the decision context. The framework extends naturally, with different computational approaches for different functional classes (linear $\rightarrow$ closed-form, general $\rightarrow$ sampling).

\subsubsection{Promising extensions}

\textbf{Other distance metrics.} The framework extends to KL divergence (likelihood shifts), $\chi^2$ divergence (specific perturbations), $f$-divergences (general family). Each captures different shift mechanisms, expanding the toolkit.

\textbf{Meta-analytic integration.} When multiple studies are available, treat each as a draw from $\cF$ and estimate functionals via empirical distributions. This bridges single-study local geometric evaluation and classical meta-analytic validation.

\textbf{Bayesian implementation.} Place a prior on the future study distribution within the local geometry and compute posterior distributions on surrogate quality $\Theta(\PP_0; \lambda)$. This provides full uncertainty quantification and enables decision-theoretic stopping rules.

\textbf{Real-time monitoring.} With 4ms inference time via closed-form concordance, bounds can be recomputed continuously as data accrue. This enables adaptive decision-making: ``Is the surrogate validated well enough to proceed?'' or ``Should we enrich for subgroups?''

\textbf{Interactive tools.} Computational efficiency enables interactive visualization of $\phi_*(\lambda)$ and worst-case scenarios. Users could upload data, select geometry, adjust $\lambda$ interactively, and explore sensitivity. This could democratize local geometric evaluation beyond specialist statisticians.

\subsection{Conclusion}

Traditional surrogate validation methods assume transportability to future studies. When this assumption is violated, inference can be misleading, with coverage dropping 20--25 percentage points below nominal levels. We introduced a framework for \textbf{local geometric evaluation}: explicitly computing worst-case surrogate performance over distributions within distance $\lambda$ of the observed data.

The framework is general, instantiated here via TV-ball DRO (arbitrary shifts) and Wasserstein DRO (covariate shifts), with sampling-based algorithms for general functionals and closed-form solutions for linear functionals. The closed-form approach provides 9--487$\times$ speedup, enabling real-time inference and large-scale sensitivity analyses.

The choice between traditional methods (assume transportability) and local geometric evaluation (evaluate transportability) depends on the decision context. For exploratory analysis with justified transportability, traditional methods are efficient. For prospective decisions with uncertain future populations, explicit evaluation provides conservative guarantees. We recommend reporting both: the conservative bound for decision-making, the descriptive estimate for context, with the gap between them quantifying transportability concern.

Software implementing the framework is available in the R package \texttt{surrogateTransportability}.

\section{Theoretical properties of the RF-ensemble approximation}

The minimax bounds $[\phi_*(\lambda), \phi^*(\lambda)]$ provide a principled answer to the question: \emph{what is the worst-case surrogate performance over all studies within total variation distance $\lambda$ of $\PP_0$?} While Section~\ref{sec:tv-implementation} established computation and finite-sample properties, we now address the fundamental question: do these bounds converge to the TV-ball minimax as $n \to \infty$?

\subsection{TV-ball minimax as the target}

For a fixed $\lambda$, define the TV-ball around $\PP_0$ as
\[
    \cB_\lambda(\PP_0) = \{\Q \in \cM_1(\Omega) : d_{\mathrm{TV}}(\Q, \PP_0) \leq \lambda\}.
\]
The \emph{TV-ball minimax} is the worst-case functional value over this ball:
\begin{align}\label{eq:tv-minimax}
    \phi_{\min}^{\mathrm{TV}}(\lambda) = \inf_{\Q \in \cB_\lambda(\PP_0)} \phi(\Q).
\end{align}
This represents an absolute worst-case over all distributions within TV distance $\lambda$ of $\PP_0$. The expected quality $\Theta(\PP_0; \lambda) = \E_{\text{Uniform}}[\phi(\Q)]$ over the uniform measure on $\cB_\lambda(\PP_0)$ satisfies $\Theta(\PP_0; \lambda) \geq \phi_{\min}^{\mathrm{TV}}(\lambda)$, providing a conservative (but less extreme) assessment of surrogate quality.

The question is: does our discretization-based method approximate $\phi_{\min}^{\mathrm{TV}}(\lambda)$?

\subsection{RF-ensemble approximation theorem}

When covariates are continuous and treatment effects $\tau_S(\bX)$, $\tau_Y(\bX)$ vary smoothly, the finite-support assumption requires discretization. Different discretization schemes explore different aspects of the treatment effect distribution. The following informal theorem characterizes convergence:

\begin{theorem}[RF-Ensemble Approximation, Informal]\label{thm:rf-ensemble}
Suppose treatment effects $\tau_S(\bX)$ and $\tau_Y(\bX)$ are Lipschitz continuous with respect to $\bX$. Let $\cJ = \{J_1, \ldots, J_K\}$ be a collection of discretization schemes with $J_k \to \infty$ as $n \to \infty$. For each scheme $k$, let $\hat{\phi}_{\min}^{(k)}(\lambda)$ be the estimated minimax from Eq.~\eqref{eq:phi-star-lower} using $J_k$ bins. Define the ensemble estimator
\[
    \hat{\phi}_{\min}^{\mathrm{ens}}(\lambda) = \min_{k=1,\ldots,K} \hat{\phi}_{\min}^{(k)}(\lambda).
\]
Under regularity conditions, as $n \to \infty$ and $J_k \to \infty$ with $J_k = o(n)$,
\[
    \hat{\phi}_{\min}^{\mathrm{ens}}(\lambda) \xrightarrow{p} \phi_{\min}^{\mathrm{TV}}(\lambda).
\]
\end{theorem}

\paragraph{Proof strategy (sketch).}
\begin{enumerate}
    \item \textbf{Random forest consistency}: RFs trained on treatment effect estimates $\hat{\tau}(\bX)$ consistently estimate the true functions $\tau(\bX)$ (Wager \& Athey, 2018). The induced partition creates regions where $\tau(\bX)$ is approximately constant.

    \item \textbf{Discretization approximation}: Any distribution $\Q \in \cB_\lambda(\PP_0)$ induces a distribution $G^\Q$ over treatment effect pairs $(\tau_S(\bX), \tau_Y(\bX))$. With $J \to \infty$, discretized distributions $G^{\Q, J}$ approximate $G^\Q$ in Wasserstein distance.

    \item \textbf{Functional continuity}: The functional $\phi$ (e.g., correlation) varies smoothly with the induced distribution $G$. Wasserstein continuity ensures $|\phi(G^{\Q, J}) - \phi(G^\Q)| = O(J^{-\alpha})$ for $\alpha > 0$ depending on smoothness.

    \item \textbf{Ensemble covering}: Using multiple discretization schemes (RF-based, quantile-based, k-means) ensures that for any adversarial $\Q^* \in \cB_\lambda$, at least one scheme produces a partition that captures the heterogeneity induced by $\Q^*$. Taking the minimum over schemes yields $\inf_k \phi(G^{\Q^*, J_k}) \to \phi(G^{\Q^*})$.

    \item \textbf{Estimation consistency}: For each scheme, the plug-in estimator $\hat{\phi}_{\min}^{(k)}$ converges to the discretized minimax. The minimum over finitely many consistent estimators is consistent.
\end{enumerate}

\subsection{Practical implications}

\paragraph{Approximation quality.}
Empirical validation confirms the theorem: with $n = 2000$ and $J \in \{9, 16, 25\}$, approximation error is consistently below $2\%$ across correlation values in $[-0.8, 0.95]$, demonstrating that the RF-ensemble with deterministic reweighting accurately approximates the TV-ball minimax.

\paragraph{Convergence rate.}
Observed error reduction from $n = 500$ ($\sim 10\%$ error) to $n = 4000$ ($\sim 1\%$ error) is consistent with $O(n^{-1/2})$ or better, suggesting the method achieves near-parametric rates despite the nonparametric discretization.

\paragraph{Number of schemes.}
Using $K = 3$--$5$ diverse schemes (e.g., RF-based, age-risk, age-biomarker, risk-biomarker, k-means) provides substantial improvement ($10$--$20\%$) over single-scheme approaches while keeping computation tractable. Additional schemes beyond $K = 5$ show diminishing returns.

\paragraph{Choice of $J$.}
For each scheme, using $J \in \{9, 16, 25\}$ and taking the minimum balances discretization error (decreases with $J$) and estimation noise (increases with $J$). Adaptive choice based on effective sample size per bin (target: $\geq 20$ observations/bin) works well in practice.

\subsection{Limitations and extensions}

\paragraph{Known limitations.}
The method requires adequate treatment effect heterogeneity: when treatment effects are near-constant across $\bX$, all future studies yield similar $(\Delta_S, \Delta_Y)$ pairs regardless of covariate distribution shifts, making surrogate evaluation trivial (but not problematic). For weak correlations ($\rho < 0.5$) with many latent subgroups ($K \geq 50$), discretization with moderate $J$ may under-represent heterogeneity, leading to conservative (wider) bounds.

\paragraph{Continuous approximation.}
An alternative to discretization is kernel-based estimation directly on continuous $\bX$. However, this requires bandwidth selection and scales poorly to high-dimensional $\bX$. The discretization approach is more robust and computationally efficient.

\paragraph{Formal proof.}
A rigorous proof of Theorem~\ref{thm:rf-ensemble} requires formalizing the covering property of the ensemble and establishing uniform convergence rates. This is left for future work but is strongly supported by the empirical validation.

% To operationalize this, we assume that $\PP_0 \sim \text{Dirichlet}(\balpha)$ for $\balpha = (\alpha_1, ..., \alpha_n), \alpha_k >0$, which results in a posterior distribution of $\cF | \PP_0 \sim \text{Dirichlet}(\balpha + 1).$ Choosing $\balpha$ encodes how similar new studies are expected to be to the current study on $\PP_0$. As $\balpha \rightarrow \infty$, new studies will be identical to the current study. 

% To better understand how choices of $\balpha$ impact the posterior distribution, consider the toy examples in Figure \ref{fig:alpha-toy-example}. When the $\balpha$ is uniform, the distribution of future studies is centered on $\PP_0$, with the magnitude of $\balpha$ determining whether future studies are mostly similar to the current study (top right in Figure \ref{fig:alpha-toy-example}, $\balpha = (5, 5, 5)$) or more diffuse (top left in the figure). When $\balpha$ is not uniform, the posterior distribution may be systematically different from the current study, primarily including studies with certain types of individuals over-represented compared to $\PP_0$ and others under-represented (bottom row of the figure). 

% \begin{figure}[htbp]
%     \centering
%     \includegraphics[width=0.8\textwidth]{dirichlet_contours_clean.png}
%     \caption{Toy examples of $\balpha$ impact on posterior distribution $\cF | \PP_0$ when $n = 3$. Contours indicate the density of the posterior, with lighter colors indicating higher density. Areas near the middle of the triangle correspond to similar studies to the current study, while areas near the edges correspond to studies with distributions that differ more from $\PP_0$.}
%     \label{fig:alpha-toy-example}
% \end{figure}

% \subsection{Procedure for a pre-specified $\balpha$}

% Depending on the scientific context, researchers may choose a specific $\balpha$ that describes their knowledge of future studies. For example, if the current study includes both young and old patients but most future studies will be focused primarily on older patients, they may choose $\balpha$ to be weighted more toward older patients. Or $\balpha$ may be chosen to be uniform such that $\alpha_k = c$ for each $k$ and some $c > 0$, where $c$ may be chosen to control the spread in the posterior distribution. 

% Let's consider some examples. We will stick with the example in \ref{tab:binary_combinations} for the moment, and we will select $\cF$ to be Dirichlet($\bpi$), where $\bpi = (\pi_{0000}, \pi_{0001}, ..., \pi_{1111})$. If we select an improper/uninformative prior on $\cF$, where the prior parameters don't contribute anything to the resulting posterior, then we obtain $\cF|\P_0 \sim \text{Dirichlet}(K\bp)$, where $\bp = (p_{0000}, ..., p_{1111})$, the observed probability masses, and $K$ is the tuning parameter controlling the spread of the distribution. 

% \subsection{Using correlation as the measure}

% If what we want is to estimate (I realize this doesn't fit exactly into the format of \eqref{posterior}, but let's worry about that later)
% \begin{align*}
%     \cor_{\cF | \P_0}\{\Delta_y(\Q), \Delta_S(\Q)\}
% \end{align*}
% then we can recover this in closed form, and the value doesn't depend on the tuning parameter $K$. In particular, we have that
% \begin{align*}
%     \cor_{\cF | \P_0}\{\Delta_y(\Q), \Delta_S(\Q)\} = \frac{\ba\trans\Sigma\bb}{\sqrt{\ba\trans\Sigma\ba \bb\trans\Sigma\bb}}
% \end{align*}
% where $\ba = (0, 1, -1, 0, 0, 1, -1, 0, 0, 1, -1, 0, 0, 1, -1, 0)$ such that $\ba\trans\bp = \Delta_Y$ and $\bb = (0, 0, 0, 0, 1, 1, 1, 1, -1, -1, -1, -1, 0, 0, 0, 0)$ such that $\bb\trans\bp = \Delta_S$, and 
% \begin{align*}
%     \Sigma = \Cov(\bp) = \inv{(K+1)} (\text{diag}(\bp) - \bp\bp\trans).
% \end{align*}

% We will put off for a moment considering other effect measures beyond correlation to consider how this can be adapted to settings on real data.

% \subsection{Extending this approach to real data}

% In real data, we don't have the potential outcomes. It might at first seem sensible to use the empirical means and the distribution $\bptilde = 1/n$ for each observation. Now we have $\Deltahat_Y = \ba\trans\bptilde$ where $\ba = Y\{A/\E(A) - (1-A)/\E(1-A)\}$ and $\Deltahat_S = \bb\trans\bp$ where $\bb = S\{A/\E(A) - (1-A)/\E(1-A)\}.$ However, computing
% \begin{align*}
%     \cor_{\cF | \P_0}\{\Deltahat_y(\Q), \Deltahat_S(\Q)\} = \frac{\ba\trans\Sigma\bb}{\sqrt{\ba\trans\Sigma\ba \bb\trans\Sigma\bb}}
% \end{align*}
% will be inflated by the correlation induced by $A$, and it doesn't immediately seem apparent how to avoid this.

% % Instead, we can acknowledge that we can't observe $\bp$ directly, but we can bound it given the data we observe. Specifically, we require that
% % \begin{align}
% %     \Omega\bp = \obar\label{constraint}
% % \end{align}
% % where $\obar$ is a vector containing values of $\Phat_n(A = a, S = s, Y = y)$, and $\Omega_{8\times16}$ is a binary constraint matrix that indicates which of the 16 counterfactual types is consistent with each of the eight observed $(A, S, Y)$ tuples and ensures that the distribution of the counterfactuals is consistent with the empirical distribution $\obar$. We may then compute $\cor_{\cF | \P_0}\{\Deltahat_y(\Q), \Deltahat_S(\Q)\}$ under the constraint \eqref{constraint}.

% Instead, we can write down the likelihood for the observed data in terms of $\bp$, which is
% \begin{align*}
%     \cL(\bp) = \prod_{s = 0}^1\prod_{y=0}^1 \left\{(1-\pi)\sum_{s' = 0}^1\sum_{y'=0}^1 p_{ss'yy'}\right\}^{n_{0sy}}\left\{\pi\sum_{s' = 0}^1\sum_{y'=0}^1 p_{s'sy'y}\right\}^{n_{1sy}}
% \end{align*}
% where $\pi= \P_0(A = 1)$ and $n_{asy} = \sum_{i=1}^n\cI\{S = s, Y = y, A = a\}$.
% Recalling that $\bp \sim$ Dirichlet($\bpi$), we put a prior on the entries in $\bpi$
% \begin{align*}
%     \pi_{ss'yy'} \sim \text{Gamma}(\alpha_{ss'yy'}, \beta_{ss'yy'}),
% \end{align*}
% and we compute the posterior distribution of $\bpi$.

  %!TEX root = main.tex

\section{Asymptotic Normality: Complete Proof}

\subsection{Setup and Notation}

\paragraph{Estimand.}
Let $\mu$ be a distribution over probability measures $\mathcal{Q}$ on the sample space $(\mathcal{O}, \mathcal{A})$. For a given distribution $\mathcal{Q}$, define treatment effects:
\begin{align}
\Delta_S(\mathcal{Q}) &= \mathbb{E}_{\mathcal{Q}}[S|A=1] - \mathbb{E}_{\mathcal{Q}}[S|A=0], \\
\Delta_Y(\mathcal{Q}) &= \mathbb{E}_{\mathcal{Q}}[Y|A=1] - \mathbb{E}_{\mathcal{Q}}[Y|A=0].
\end{align}

Let $\phi: \mathbb{R}^2 \to \mathbb{R}$ be a functional. The estimand of interest is:
\begin{equation}
\Theta(\mu) = \mathbb{E}_{\mu}[\phi(\Delta_S(\mathcal{Q}), \Delta_Y(\mathcal{Q}))].
\end{equation}

\paragraph{Observed data.}
We observe $n$ i.i.d.\ observations $O_1, \ldots, O_n \sim \mathbb{P}_0$, where $O_i = (X_i, A_i, S_i, Y_i)$.

\paragraph{Sampling scheme.}
We draw $\mathcal{Q}_1, \ldots, \mathcal{Q}_M$ from $\mu$ via MCMC (hit-and-run sampling).

\paragraph{Estimator.}
For each sampled $\mathcal{Q}_m$, we estimate treatment effects using reweighted data:
\begin{align}
\hat{\Delta}_S^n(\mathcal{Q}_m) &= \frac{\sum_{i=1}^n w_i(\mathcal{Q}_m) A_i S_i}{\sum_{i=1}^n w_i(\mathcal{Q}_m) A_i} - \frac{\sum_{i=1}^n w_i(\mathcal{Q}_m) (1-A_i) S_i}{\sum_{i=1}^n w_i(\mathcal{Q}_m) (1-A_i)}, \\
\hat{\Delta}_Y^n(\mathcal{Q}_m) &= \frac{\sum_{i=1}^n w_i(\mathcal{Q}_m) A_i Y_i}{\sum_{i=1}^n w_i(\mathcal{Q}_m) A_i} - \frac{\sum_{i=1}^n w_i(\mathcal{Q}_m) (1-A_i) Y_i}{\sum_{i=1}^n w_i(\mathcal{Q}_m) (1-A_i)},
\end{align}
where $w_i(\mathcal{Q}) = \frac{d\mathcal{Q}}{d\mathbb{P}_0}(X_i)$ are importance weights.

The plug-in estimator is:
\begin{equation}
\hat{\Theta}_n = \frac{1}{M} \sum_{m=1}^M \phi(\hat{\Delta}_S^n(\mathcal{Q}_m), \hat{\Delta}_Y^n(\mathcal{Q}_m)).
\end{equation}

\subsection{Assumptions}

\begin{assumption}[Treatment effect estimation]
\label{assm:treatment-effects}
For each $\mathcal{Q} \in \text{supp}(\mu)$:
\begin{enumerate}[(a)]
\item Randomization or conditional exchangeability: $Y(a), S(a) \perp A \mid X$ for $a \in \{0,1\}$ under $\mathcal{Q}$.
\item Positivity: $\mathcal{Q}(A=a|X) > 0$ a.s.\ for $a \in \{0,1\}$.
\item Bounded weights: $\sup_{\mathcal{Q} \in \text{supp}(\mu)} \sup_{x \in \text{supp}(\mathbb{P}_0)} \frac{d\mathcal{Q}}{d\mathbb{P}_0}(x) \leq C < \infty$.
\item Bounded second moments: $\mathbb{E}_{\mathbb{P}_0}[S^2 + Y^2] < \infty$.
\end{enumerate}
\end{assumption}

\begin{assumption}[Functional regularity]
\label{assm:functional}
The functional $\phi: \mathbb{R}^2 \to \mathbb{R}$ is Hadamard differentiable at $(\Delta_S(\mathcal{Q}), \Delta_Y(\mathcal{Q}))$ for all $\mathcal{Q} \in \text{supp}(\mu)$, with gradient $\nabla\phi(\Delta_S, \Delta_Y) = (\partial_1 \phi, \partial_2 \phi)$ satisfying:
\begin{equation}
\sup_{\mathcal{Q} \in \text{supp}(\mu)} \|\nabla\phi(\Delta_S(\mathcal{Q}), \Delta_Y(\mathcal{Q}))\| \leq C_\phi < \infty.
\end{equation}
\end{assumption}

\begin{remark}[Aggregate functionals]
\label{rem:aggregate-gradient}
For aggregate functionals such as correlation, $\Theta(\mu) = \varphi(m(\mu))$ is a smooth function of a vector of $\mu$-moments $m(\mu)$ (five moments for correlation; see Section~\ref{sec:correlation-application}), not a pointwise map of a single pair. Throughout, $\nabla\phi(\Delta_S(\mathcal{Q}), \Delta_Y(\mathcal{Q}))$ denotes the \emph{composed} per-study gradient $\nabla\varphi(m(\mu)) \cdot \partial m/\partial(\Delta_S(\mathcal{Q}), \Delta_Y(\mathcal{Q}))$ obtained by the chain rule; the influence-function expressions \eqref{eq:aipw-s}--\eqref{eq:aipw-y} and Lemma~\ref{lem:functional-expansion} should be read with this interpretation. Non-degeneracy (Assumption~\ref{assm:nondegeneracy}) guarantees the bounded-gradient condition above for the correlation functional.
\end{remark}

\begin{assumption}[Nuisance parameter convergence]
\label{assm:nuisance}
For observational studies with estimated propensity scores $\hat{e}(X) = P(A=1|X)$ and outcome regressions $\hat{\mu}_a(X) = \mathbb{E}[Y|A=a,X]$ for $a \in \{0,1\}$, obtained via cross-fitting:
\begin{enumerate}[(a)]
\item \textbf{Strong positivity:} $e(x) \in [c, 1-c]$ uniformly for some $c \geq 0.1$.
\item \textbf{Propensity score rate:} $\|\hat{e} - e\|_{L^2(\mathbb{P}_0)} = o_P(n^{-1/4})$.
\item \textbf{Outcome regression rate:} $\|\hat{\mu}_a - \mu_a\|_{L^2(\mathbb{P}_0)} = o_P(n^{-1/4})$ for $a \in \{0,1\}$.
\item \textbf{Product convergence:} $\|\hat{e} - e\|_{L^2} \cdot \|\hat{\mu}_a - \mu_a\|_{L^2} = o_P(n^{-1/2})$.
\item \textbf{Bounded conditional moments:} $\mathbb{E}_{\mathbb{P}_0}[Y^2 | X, A=a] \leq M < \infty$ and $\mathbb{E}_{\mathbb{P}_0}[S^2 | X, A=a] \leq M < \infty$ uniformly.
\item \textbf{Weights independent of nuisances:} $w(X; \mathcal{Q}) = \frac{d\mathcal{Q}}{d\mathbb{P}_0}(X)$ does not depend on $\eta = (e, \mu_0, \mu_1)$.
\end{enumerate}
\end{assumption}

\begin{remark}[When $o_P(n^{-1/4})$ rates hold]
The $o_P(n^{-1/4})$ rates in Assumption~\ref{assm:nuisance}(b-c) hold under standard nonparametric conditions such as:
\begin{itemize}
\item Hölder smoothness $\beta > d/2$ for the nuisance functions (satisfied by GAMs, kernel methods, neural networks with appropriate architecture)
\item Sparsity or approximate sparsity assumptions (satisfied by elastic net, boosting, adaptive LASSO)
\item Structural assumptions enabling fast rates (satisfied by random forests with appropriate tuning, gradient boosting)
\end{itemize}
See Chernozhukov et al.\ (2018) for specific function class requirements. In practice, modern machine learning methods (GAMs, random forests, gradient boosting) typically achieve these rates when properly cross-validated.
\end{remark}

\begin{assumption}[Bounded second derivatives]
\label{assm:second-derivatives}
The AIPW influence function has bounded second derivatives uniformly over nuisance parameters in a neighborhood of $\eta_0 = (e, \mu_0, \mu_1)$:
\begin{equation}
\mathbb{E}_{\mathbb{P}_0}\left[\sup_{\|\eta - \eta_0\| \leq \delta} \left\|\frac{\partial^2}{\partial\eta^2}\psi_S^{\text{AIPW}}(O; \eta)\right\|^2\right] \leq B < \infty
\end{equation}
for some $\delta > 0$, and similarly for $\psi_Y^{\text{AIPW}}$.
\end{assumption}

\begin{assumption}[Uniformly bounded IF variance under reweighting]
\label{assm:bounded-if-variance}
The influence functions have uniformly bounded second moments over $\mathcal{Q} \in \text{supp}(\mu)$:
\begin{equation}
\sup_{\mathcal{Q} \in \text{supp}(\mu)} \mathbb{E}_{\mathbb{P}_0}[(w(X; \mathcal{Q}) \psi_S^{\text{AIPW}}(O; \eta_0))^2] \leq V < \infty,
\end{equation}
and similarly for $\psi_Y^{\text{AIPW}}$. This ensures the CLT applies uniformly as $M \to \infty$.
\end{assumption}

\begin{assumption}[Lipschitz gradient of functional]
\label{assm:lipschitz}
The functional $\phi: \mathbb{R}^2 \to \mathbb{R}$ has Lipschitz continuous gradient:
\begin{equation}
\|\nabla\phi(\theta_1) - \nabla\phi(\theta_2)\| \leq L\|\theta_1 - \theta_2\|
\end{equation}
for some $L < \infty$ and all $\theta_1, \theta_2$ in the range of $(\Delta_S(\mathcal{Q}), \Delta_Y(\mathcal{Q}))$ for $\mathcal{Q} \in \text{supp}(\mu)$.
\end{assumption}

\begin{assumption}[MCMC convergence]
\label{assm:mcmc}
The hit-and-run sampler satisfies geometric ergodicity with stationary distribution $\mu$. After burn-in, the empirical measure satisfies:
\begin{equation}
\left\|\hat{\mu}_M - \mu\right\|_{\text{TV}} = O_P(M^{-1/2}),
\end{equation}
where $\hat{\mu}_M = \frac{1}{M}\sum_{m=1}^M \delta_{\mathcal{Q}_m}$ and $\|\cdot\|_{\text{TV}}$ denotes total variation distance.
\end{assumption}

\begin{assumption}[Geometry regularity]
\label{assm:geometry}
The distribution $\mu$ has compact support in the weak topology: $\text{supp}(\mu)$ is compact.
\end{assumption}

\begin{assumption}[Sample size and MCMC regime]
\label{assm:sample-size}
$n \to \infty$. The primary (conditional) result treats the MCMC samples $\{\mathcal{Q}_m\}_{m=1}^M$ as fixed and holds for any $M \geq M_0$ large enough for the sampler to have reached its stationary regime. The \emph{unconditional} result, centered at the population target $\Theta(\mu)$, additionally requires $M = M(n) \to \infty$ with $n/M \to 0$ (equivalently $n = o(M)$), so that the MCMC term is asymptotically negligible relative to the $\sqrt{n}$ estimation term.
\end{assumption}

\begin{assumption}[Non-degeneracy of across-study variances]
\label{assm:nondegeneracy}
There exists $c_0 > 0$ such that $\mathrm{Var}_\mu(\Delta_S(\mathcal{Q})) \geq c_0$ and $\mathrm{Var}_\mu(\Delta_Y(\mathcal{Q})) \geq c_0$. This guarantees the correlation functional and its gradient are well-defined; the asymptotic variance scales like $c_0^{-3}$ as the boundary is approached.
\end{assumption}

\begin{remark}[MCMC sample size and the two regimes]
Decompose $\sqrt{n}(\hat{\Theta}_n - \Theta(\mu)) = \sqrt{n}(\hat{\Theta}_n - \Theta(\hat{\mu}_M)) + \sqrt{n}(\Theta(\hat{\mu}_M) - \Theta(\mu))$. The first term is $O_P(1)$ (the conditional CLT); the second is $\sqrt{n}\cdot O_P(M^{-1/2}) = O_P((n/M)^{1/2})$ by Lemma~\ref{lem:mcmc-error}. Hence:
\begin{itemize}
\item \textbf{Conditional regime (any moderate $M$):} inference targets the $M$-study correlation $\Theta(\hat{\mu}_M)$; the estimator's reported standard error is the conditional SE and is valid for $\Theta(\hat{\mu}_M)$. The residual $\Theta(\hat{\mu}_M) - \Theta(\mu) = O_P(M^{-1/2})$ is an MCMC approximation bias that does not vanish.
\item \textbf{Unconditional regime ($n = o(M)$):} the MCMC term $O_P((n/M)^{1/2}) = o_P(1)$ vanishes and $\sqrt{n}(\hat{\Theta}_n - \Theta(\mu)) \xrightarrow{d} N(0,\sigma^2(\mu))$.
\end{itemize}
Note this requires $M$ to grow \emph{faster} than $n$ (not $M = o(n)$). In the numerical studies $M \approx 2100 \gg n^{1/2}$ per effective cell, and the estimation term dominates; we report the conditional SE.
\end{remark}

\subsection{Main Result}

\begin{theorem}[Asymptotic normality, conditional on MCMC samples]
\label{thm:asymptotic-normality}
Under Assumptions~\ref{assm:treatment-effects}--\ref{assm:sample-size}, conditional on the MCMC samples $\{\mathcal{Q}_1, \ldots, \mathcal{Q}_M\}$, as $n \to \infty$:
\begin{equation}
\sqrt{n}(\hat{\Theta}_n - \Theta(\hat{\mu}_M)) \xrightarrow{d} N(0, \sigma^2(\hat{\mu}_M)),
\end{equation}
where $\Theta(\hat{\mu}_M) = \mathbb{E}_{\hat{\mu}_M}[\phi(\Delta_S(\mathcal{Q}), \Delta_Y(\mathcal{Q}))]$ and the asymptotic variance is:
\begin{equation}
\sigma^2(\hat{\mu}_M) = \mathbb{E}_{\mathbb{P}_0}[\psi_\Theta(O; \hat{\mu}_M)^2],
\end{equation}
with influence function:
\begin{equation}
\psi_\Theta(O; \hat{\mu}_M) = \mathbb{E}_{\hat{\mu}_M}\left[\nabla\phi(\Delta_S(\mathcal{Q}), \Delta_Y(\mathcal{Q})) \cdot \begin{pmatrix} \psi_S^{\text{AIPW}}(O; \eta_0, \mathcal{Q}) \\ \psi_Y^{\text{AIPW}}(O; \eta_0, \mathcal{Q}) \end{pmatrix}\right],
\end{equation}
where $\psi_S^{\text{AIPW}}(O; \eta, \mathcal{Q})$ and $\psi_Y^{\text{AIPW}}(O; \eta, \mathcal{Q})$ are the AIPW influence functions defined in \eqref{eq:aipw-s}--\eqref{eq:aipw-y}. For randomized trials, these reduce to weighted difference-in-means; for observational studies with cross-fitting, they provide doubly robust estimation.

Furthermore:
\begin{enumerate}[(a)]
\item If $M$ is fixed (or grows slower than $n$), $\Theta(\hat{\mu}_M) - \Theta(\mu) = O_P(M^{-1/2})$ (MCMC approximation bias), and inference is for $\Theta(\hat{\mu}_M)$.
\item If $M = M(n) \to \infty$ with $n/M \to 0$, the MCMC term $\sqrt{n}(\Theta(\hat{\mu}_M) - \Theta(\mu)) = O_P((n/M)^{1/2}) = o_P(1)$, and unconditionally:
\begin{equation}
\sqrt{n}(\hat{\Theta}_n - \Theta(\mu)) \xrightarrow{d} N(0, \sigma^2(\mu)).
\end{equation}
\end{enumerate}
\end{theorem}

\begin{remark}[Practical interpretation]
Conditioning on MCMC samples separates two sources of randomness:
\begin{itemize}
\item \textbf{Sampling variation} (from finite $n$): Captured by the asymptotic distribution, rate $n^{-1/2}$.
\item \textbf{MCMC approximation} (from finite $M$): Bias in target, magnitude $M^{-1/2}$.
\end{itemize}
When $M$ is large relative to $n$, the MCMC bias $O_P(M^{-1/2})$ is negligible relative to the $O_P(n^{-1/2})$ sampling variation; formally it vanishes in the regime $n = o(M)$.
\end{remark}

\subsection{Proof of Theorem~\ref{thm:asymptotic-normality}}

\subsubsection{Proof roadmap}

The proof proceeds in four main steps:
\begin{enumerate}
\item \textbf{MCMC approximation:} Show (via a Markov chain CLT) that the empirical measure over MCMC samples approximates $\mu$ at rate $M^{-1/2}$. Relative to the $\sqrt{n}$ estimation term this contributes $O_P((n/M)^{1/2})$, negligible only in the unconditional regime $n = o(M)$ (Assumption~\ref{assm:sample-size}); otherwise inference is conditional on $\hat{\mu}_M$.
\item \textbf{Treatment effect expansion (AIPW):} Establish Neyman orthogonality under reweighting (Lemma~\ref{lem:orthogonality}), verify conditions for cross-fitting with bounded weights (Lemma~\ref{lem:cross-fitting}), and obtain uniform influence function expansion with $o_P(n^{-1/2})$ remainder. Key innovation: bounded second derivatives + bounded weights give deterministic uniformity over MCMC samples.
\item \textbf{Functional delta method:} Apply Hadamard differentiability to obtain the expansion for $\phi(\hat{\Delta}_S^n(\mathcal{Q}), \hat{\Delta}_Y^n(\mathcal{Q}))$.
\item \textbf{Integration and CLT:} Average over $\mu$ (or its MCMC approximation) and apply the central limit theorem using uniform moment bounds.
\end{enumerate}

\subsubsection{Step 1: MCMC approximation}

Let $\hat{\mu}_M = \frac{1}{M}\sum_{m=1}^M \delta_{\mathcal{Q}_m}$ denote the empirical measure over MCMC samples.

\begin{lemma}[MCMC approximation error]
\label{lem:mcmc-error}
Under Assumptions~\ref{assm:mcmc} (geometric ergodicity) and \ref{assm:geometry} (compact support), for any bounded measurable function $g: \text{supp}(\mu) \to \mathbb{R}$ with $\|g\|_\infty \leq 1$:
\begin{equation}
\left|\mathbb{E}_{\hat{\mu}_M}[g(\mathcal{Q})] - \mathbb{E}_{\mu}[g(\mathcal{Q})]\right| = O_P(M^{-1/2}).
\end{equation}
\end{lemma}

\begin{proof}[Proof of Lemma~\ref{lem:mcmc-error}]
The hit-and-run samples $\mathcal{Q}_1, \ldots, \mathcal{Q}_M$ form a \emph{dependent} Markov chain, not independent draws; we therefore invoke the Markov chain CLT rather than the i.i.d.\ CLT. By Assumption~\ref{assm:mcmc}, the chain is geometrically ergodic with stationary distribution $\mu$, and $g$ is bounded (hence in $L^2(\mu)$ with all polynomial moments). Under geometric ergodicity and a bounded (thus $L^{2+\delta}(\mu)$) functional, the Markov chain CLT holds \cite{jones2004markov, roberts2004general}:
\begin{equation}
\sqrt{M}\left(\frac{1}{M}\sum_{m=1}^M g(\mathcal{Q}_m) - \mathbb{E}_{\mu}[g]\right) \xrightarrow{d} N\!\left(0, \sigma_g^2\right),
\end{equation}
where $\sigma_g^2 = \mathrm{Var}_\mu(g(\mathcal{Q}_1)) + 2\sum_{k=1}^\infty \mathrm{Cov}_\mu(g(\mathcal{Q}_1), g(\mathcal{Q}_{1+k}))$ is the \emph{integrated-autocorrelation} asymptotic variance (finite by geometric ergodicity). Hence $\mathbb{E}_{\hat{\mu}_M}[g] - \mathbb{E}_{\mu}[g] = O_P(M^{-1/2})$. The rate is unchanged from the independent case; only the variance constant differs (it inflates by the integrated autocorrelation time). Applying this to $g = \phi(\Delta_S(\cdot), \Delta_Y(\cdot))$ (bounded under Assumptions~\ref{assm:treatment-effects} and \ref{assm:functional}) gives $\Theta(\hat{\mu}_M) - \Theta(\mu) = O_P(M^{-1/2})$.
\end{proof}

\begin{corollary}
\label{cor:mcmc-negligible}
Under Assumption~\ref{assm:sample-size}:
\begin{enumerate}[(a)]
\item If $M$ is fixed, $|\mathbb{E}_{\hat{\mu}_M}[g] - \mathbb{E}_{\mu}[g]| = O_P(M^{-1/2})$ is constant in $n$.
\item If $M \to \infty$, $|\mathbb{E}_{\hat{\mu}_M}[g] - \mathbb{E}_{\mu}[g]| = o_P(1)$ vanishes as $n \to \infty$.
\end{enumerate}
\end{corollary}

\begin{remark}
In the fixed-$M$ case, the theorem establishes asymptotic normality for $\Theta(\hat{\mu}_M)$ (the MCMC-approximated target) rather than $\Theta(\mu)$ itself. The difference $|\Theta(\hat{\mu}_M) - \Theta(\mu)| = O_P(M^{-1/2})$ is a source of bias that does not vanish unless $M \to \infty$. It is negligible relative to $\sqrt{n}$ sampling variation precisely in the regime $n = o(M)$.
\end{remark}

\subsubsection{Step 2: Treatment effect expansion}

\paragraph{AIPW influence functions.}
For observational studies with estimated nuisances, we use augmented inverse probability weighting (AIPW). The influence functions are:
\begin{align}
\psi_S^{\text{AIPW}}(O; \eta, \mathcal{Q}) &= w(X; \mathcal{Q}) \Bigg[\frac{A(S - \mu_1^S(X))}{e(X)} - \frac{(1-A)(S - \mu_0^S(X))}{1-e(X)} \nonumber \\
&\quad + \mu_1^S(X) - \mu_0^S(X)\Bigg] - \Delta_S(\mathcal{Q}), \label{eq:aipw-s} \\
\psi_Y^{\text{AIPW}}(O; \eta, \mathcal{Q}) &= w(X; \mathcal{Q}) \Bigg[\frac{A(Y - \mu_1^Y(X))}{e(X)} - \frac{(1-A)(Y - \mu_0^Y(X))}{1-e(X)} \nonumber \\
&\quad + \mu_1^Y(X) - \mu_0^Y(X)\Bigg] - \Delta_Y(\mathcal{Q}), \label{eq:aipw-y}
\end{align}
where:
\begin{itemize}
\item $\eta = (e, \mu_0^S, \mu_1^S, \mu_0^Y, \mu_1^Y)$ are nuisance parameters,
\item $e(X) = P(A=1|X)$ is the propensity score,
\item $\mu_a^S(X) = \mathbb{E}[S|A=a,X]$ and $\mu_a^Y(X) = \mathbb{E}[Y|A=a,X]$ are outcome regressions,
\item $w(X; \mathcal{Q}) = \frac{d\mathcal{Q}}{d\mathbb{P}_0}(X)$ are importance weights (independent of $\eta$ by Assumption~\ref{assm:nuisance}(f)),
\item $\Delta_S(\mathcal{Q}) = \mathbb{E}_{\mathcal{Q}}[\mu_1^S(X) - \mu_0^S(X)]$ and $\Delta_Y(\mathcal{Q}) = \mathbb{E}_{\mathcal{Q}}[\mu_1^Y(X) - \mu_0^Y(X)]$.
\end{itemize}

\begin{remark}[Three-term decomposition]
The AIPW influence function has three terms:
\begin{enumerate}
\item \textbf{IPW term:} $\frac{A(S - \mu_1^S(X))}{e(X)} - \frac{(1-A)(S - \mu_0^S(X))}{1-e(X)}$ (inverse probability weighting of residuals)
\item \textbf{Regression term:} $\mu_1^S(X) - \mu_0^S(X)$ (predicted treatment effect)
\item \textbf{Centering:} the \emph{constant} $-\Delta_S(\mathcal{Q})$, subtracted \emph{outside} the importance weight, ensuring $\mathbb{E}_{\mathbb{P}_0}[\psi_S^{\text{AIPW}}] = 0$.
\end{enumerate}
This structure provides \textbf{double robustness}: the estimator is consistent if either the propensity score $e(X)$ or the outcome regressions $\mu_a(X)$ are correctly specified (not both required).

Two features are specific to this estimand and distinguish it from the textbook difference-in-means AIPW score. First, the entire bracketed AIPW score, \emph{including} the regression term $\mu_1^S(X) - \mu_0^S(X)$, is multiplied by the importance weight $w(X;\mathcal{Q})$. This is correct because the target $\Delta_S(\mathcal{Q}) = \mathbb{E}_{\mathcal{Q}}[\mu_1^S(X) - \mu_0^S(X)]$ is an expectation under the \emph{fixed, external} composition $\mathcal{Q}$, not under $\mathbb{P}_0$; there is consequently no $\mathbb{P}_0$-covariate-law augmentation term. Second, the centering is the constant $-\Delta_S(\mathcal{Q})$, not $-w(X;\mathcal{Q})\Delta_S(\mathcal{Q})$: only constant centering yields a mean-zero per-observation influence function with the correct variance (see the companion derivation, \texttt{derivation\_influence\_functions.md}).
\end{remark}

\begin{remark}[Special case: RCTs]
For randomized controlled trials where $A \perp (S, Y, X)$ under $\mathbb{P}_0$, the propensity score $e(X) = P(A=1)$ is known and constant. Setting $\mu_a^S(X) = \mathbb{E}[S|A=a]$ and $\mu_a^Y(X) = \mathbb{E}[Y|A=a]$ (constants), the AIPW IF reduces to the simple weighted difference-in-means IF used in the RCT-only proof. Thus, the AIPW framework generalizes the RCT case.
\end{remark}

\begin{lemma}[Neyman orthogonality under reweighting]
\label{lem:orthogonality}
Under Assumptions~\ref{assm:treatment-effects} and \ref{assm:nuisance}, for the weighted AIPW functional $\theta_w(\eta; \mathcal{Q}) = \mathbb{E}_{\mathbb{P}_0}[w(X; \mathcal{Q}) \psi_S^{\text{AIPW}}(O; \eta, \mathcal{Q})]$:
\begin{equation}
\frac{\partial}{\partial\eta}\theta_w(\eta; \mathcal{Q})\bigg|_{\eta = \eta_0}[r] = 0
\end{equation}
for all $r$ in the tangent space and all $\mathcal{Q} \in \text{supp}(\mu)$.
\end{lemma}

\begin{proof}[Proof of Lemma~\ref{lem:orthogonality}]
We prove the result for $\psi_S^{\text{AIPW}}$; the proof for $\psi_Y^{\text{AIPW}}$ is identical.

\paragraph{Step 1: Decompose AIPW score.}
Write the AIPW score as:
\begin{equation}
\psi_S^{\text{AIPW}}(O; \eta, \mathcal{Q}) = \psi_S^{\text{IPW}}(O; \eta) + \psi_S^{\text{reg}}(O; \eta) - \Delta_S(\mathcal{Q}),
\end{equation}
where:
\begin{align}
\psi_S^{\text{IPW}}(O; \eta) &= \frac{A(S - \mu_1^S(X))}{e(X)} - \frac{(1-A)(S - \mu_0^S(X))}{1-e(X)}, \\
\psi_S^{\text{reg}}(O; \eta) &= \mu_1^S(X) - \mu_0^S(X).
\end{align}

\paragraph{Step 2: Standard AIPW orthogonality (conditional).}
For standard AIPW (without reweighting), Neyman orthogonality is well-established (Chernozhukov et al., 2018): the conditional expectation of the pathwise derivative vanishes:
\begin{equation}
\mathbb{E}\left[\frac{\partial}{\partial\eta}\psi_S^{\text{AIPW}}(O; \eta_0)[r] \,\bigg|\, X\right] = 0
\end{equation}
for all $r$ in the tangent space. This is the \textbf{conditional mean-zero property}.

\paragraph{Step 3: Weights don't depend on nuisances.}
By Assumption~\ref{assm:nuisance}(f), $w(X; \mathcal{Q}) = \frac{d\mathcal{Q}}{d\mathbb{P}_0}(X)$ is a function of $X$ only and does not depend on $\eta = (e, \mu_0, \mu_1)$.

\paragraph{Step 4: Orthogonality survives weighting.}
For the weighted functional:
\begin{align}
\frac{\partial}{\partial\eta}\theta_w(\eta_0; \mathcal{Q})[r] &= \frac{\partial}{\partial\eta}\mathbb{E}_{\mathbb{P}_0}\left[w(X; \mathcal{Q}) \psi_S^{\text{AIPW}}(O; \eta_0, \mathcal{Q})\right][r] \\
&= \mathbb{E}_{\mathbb{P}_0}\left[w(X; \mathcal{Q}) \cdot \frac{\partial}{\partial\eta}\psi_S^{\text{AIPW}}(O; \eta_0, \mathcal{Q})[r]\right] \\
&= \mathbb{E}_{\mathbb{P}_0}\left[w(X; \mathcal{Q}) \cdot \mathbb{E}\left[\frac{\partial}{\partial\eta}\psi_S^{\text{AIPW}}(O; \eta_0)[r] \,\bigg|\, X\right]\right] \\
&= \mathbb{E}_{\mathbb{P}_0}[w(X; \mathcal{Q}) \cdot 0] = 0.
\end{align}
The third equality uses the law of iterated expectations, and the fourth uses the conditional mean-zero property from Step 2.

\paragraph{Conclusion.}
Neyman orthogonality holds for the weighted functional. This is the key property that eliminates first-order nuisance bias.
\end{proof}

\begin{lemma}[Cross-fitting with bounded weights]
\label{lem:cross-fitting}
Under Assumptions~\ref{assm:treatment-effects}, \ref{assm:nuisance}, and \ref{assm:second-derivatives}, the cross-fitted AIPW estimator with bounded weights satisfies:
\begin{equation}
\hat{\Delta}_S(\mathcal{Q}) - \Delta_S(\mathcal{Q}) = \frac{1}{n}\sum_{i=1}^n w_i(\mathcal{Q}) \psi_S^{\text{AIPW}}(O_i; \eta_0, \mathcal{Q}) + o_P(n^{-1/2}),
\end{equation}
uniformly over $\mathcal{Q} \in \text{supp}(\mu)$. Similarly for $\hat{\Delta}_Y(\mathcal{Q})$.
\end{lemma}

\begin{proof}[Proof of Lemma~\ref{lem:cross-fitting}]
We verify the conditions of Theorem 3.1 in Chernozhukov et al.\ (2018) for the weighted AIPW estimator.

\paragraph{Condition 1: Neyman orthogonality.}
Verified in Lemma~\ref{lem:orthogonality}. The weighted functional $\theta_w(\eta; \mathcal{Q})$ satisfies $\frac{\partial}{\partial\eta}\theta_w(\eta_0; \mathcal{Q}) = 0$.

\paragraph{Condition 2: $L^2$ convergence rates.}
By Assumption~\ref{assm:nuisance}(b-c):
\begin{align}
\|\hat{e} - e\|_{L^2(\mathbb{P}_0)} &= o_P(n^{-1/4}), \\
\|\hat{\mu}_a - \mu_a\|_{L^2(\mathbb{P}_0)} &= o_P(n^{-1/4}) \quad \text{for } a \in \{0,1\}.
\end{align}

\paragraph{Condition 3: Bounded moments.}
By Assumption~\ref{assm:nuisance}(e), $\mathbb{E}[Y^2|X,A], \mathbb{E}[S^2|X,A] \leq M < \infty$, ensuring the influence function has finite second moments.

\paragraph{Condition 4: Cross-fitting.}
Sample splitting into $K$ folds ensures nuisance estimators $\hat{\eta}^{(-k)}$ are independent of observations in fold $k$, eliminating overfitting bias.

\paragraph{Condition 5: Bounded weights preserve all properties.}
By Assumption~\ref{assm:treatment-effects}(c), $w_i(\mathcal{Q}) \leq C < \infty$ uniformly. Bounded weights preserve:
\begin{itemize}
\item Neyman orthogonality (Lemma~\ref{lem:orthogonality})
\item $L^2$ convergence of weighted IF: $\|w \cdot \psi^{\text{AIPW}}(\cdot; \hat{\eta}) - w \cdot \psi^{\text{AIPW}}(\cdot; \eta_0)\|_{L^2} \leq C \|\psi^{\text{AIPW}}(\cdot; \hat{\eta}) - \psi^{\text{AIPW}}(\cdot; \eta_0)\|_{L^2}$
\item Bounded second moments (Assumption~\ref{assm:bounded-if-variance})
\end{itemize}

\paragraph{Second-order remainder bound.}
By Assumption~\ref{assm:second-derivatives}, second derivatives of $\psi_S^{\text{AIPW}}$ are bounded. Taylor expansion gives:
\begin{align}
&\hat{\Delta}_S(\mathcal{Q}) - \Delta_S(\mathcal{Q}) \\
&= \frac{1}{n}\sum_{i=1}^n w_i(\mathcal{Q}) \psi_S^{\text{AIPW}}(O_i; \eta_0, \mathcal{Q}) + \frac{1}{n}\sum_{i=1}^n w_i(\mathcal{Q}) R_i(\hat{\eta}, \eta_0),
\end{align}
where the remainder satisfies:
\begin{equation}
R_i(\hat{\eta}, \eta_0) = \int_0^1 (1-t) \frac{\partial^2}{\partial\eta^2}\psi_S^{\text{AIPW}}(O_i; \tilde{\eta}_t, \mathcal{Q})[\hat{\eta} - \eta_0]^2 dt,
\end{equation}
with $\tilde{\eta}_t = \eta_0 + t(\hat{\eta} - \eta_0)$.

By Assumption~\ref{assm:second-derivatives} and bounded weights:
\begin{equation}
\left|\frac{1}{n}\sum_{i=1}^n w_i(\mathcal{Q}) R_i\right| \leq C \cdot B \cdot \|\hat{\eta} - \eta_0\|_{L^2}^2 = C \cdot B \cdot o_P(n^{-1/2}),
\end{equation}
where the last equality uses Assumption~\ref{assm:nuisance}(d).

\paragraph{Uniformity over $\mathcal{Q}$.}
The remainder bound $|R_n(\mathcal{Q})| \leq CB\|\hat{\eta} - \eta_0\|_{L^2}^2$ is \textbf{independent of $\mathcal{Q}$} (depends only on $C, B, \|\hat{\eta} - \eta_0\|_{L^2}$). Therefore:
\begin{equation}
\sup_{\mathcal{Q} \in \text{supp}(\mu)} |R_n(\mathcal{Q})| = o_P(n^{-1/2}).
\end{equation}
This is \textbf{deterministic uniformity}—no maximal inequalities or Donsker theory required.

\paragraph{Conclusion.}
By Chernozhukov et al.\ (2018, Theorem 3.1), the expansion holds with $o_P(n^{-1/2})$ remainder, uniformly over $\mathcal{Q}$.
\end{proof}


\subsubsection{Step 3: Functional delta method}

\begin{lemma}[Functional expansion]
\label{lem:functional-expansion}
Under Assumptions~\ref{assm:treatment-effects}, \ref{assm:nuisance}, \ref{assm:functional}, and \ref{assm:lipschitz}, uniformly over $\mathcal{Q} \in \text{supp}(\mu)$:
\begin{align}
&\sqrt{n}(\phi(\hat{\Delta}_S^n(\mathcal{Q}), \hat{\Delta}_Y^n(\mathcal{Q})) - \phi(\Delta_S(\mathcal{Q}), \Delta_Y(\mathcal{Q}))) \nonumber \\
&\quad = \nabla\phi(\Delta_S(\mathcal{Q}), \Delta_Y(\mathcal{Q})) \cdot \begin{pmatrix} \frac{1}{\sqrt{n}}\sum_{i=1}^n \psi_S^{\text{AIPW}}(O_i; \eta_0, \mathcal{Q}) \\ \frac{1}{\sqrt{n}}\sum_{i=1}^n \psi_Y^{\text{AIPW}}(O_i; \eta_0, \mathcal{Q}) \end{pmatrix} + o_P(1).
\end{align}
\end{lemma}

\begin{proof}[Proof of Lemma~\ref{lem:functional-expansion}]
By Hadamard differentiability of $\phi$ (Assumption~\ref{assm:functional}):
\begin{align}
&\phi(\hat{\Delta}_S^n(\mathcal{Q}), \hat{\Delta}_Y^n(\mathcal{Q})) - \phi(\Delta_S(\mathcal{Q}), \Delta_Y(\mathcal{Q})) \\
&\quad = \nabla\phi(\Delta_S(\mathcal{Q}), \Delta_Y(\mathcal{Q})) \cdot \begin{pmatrix} \hat{\Delta}_S^n(\mathcal{Q}) - \Delta_S(\mathcal{Q}) \\ \hat{\Delta}_Y^n(\mathcal{Q}) - \Delta_Y(\mathcal{Q}) \end{pmatrix} + o(\|(\hat{\Delta}_S^n, \hat{\Delta}_Y^n) - (\Delta_S, \Delta_Y)\|).
\end{align}

By Lemma~\ref{lem:cross-fitting}, $\|(\hat{\Delta}_S^n, \hat{\Delta}_Y^n) - (\Delta_S, \Delta_Y)\| = O_P(n^{-1/2})$, hence the remainder is $o_P(n^{-1/2})$. Substituting the expansion from Lemma~\ref{lem:cross-fitting} gives the result.
\end{proof}

\subsubsection{Step 4: Integration and CLT}

\begin{proof}[Proof of Theorem~\ref{thm:asymptotic-normality}]
We combine the preceding lemmas in four steps.

\paragraph{Step 4.1: Condition on MCMC samples.}
Throughout this proof, we condition on the MCMC samples $\{\mathcal{Q}_1, \ldots, \mathcal{Q}_M\}$, treating $\hat{\mu}_M$ as fixed. The estimator becomes:
\begin{equation}
\hat{\Theta}_n = \mathbb{E}_{\hat{\mu}_M}[\phi(\hat{\Delta}_S^n(\mathcal{Q}), \hat{\Delta}_Y^n(\mathcal{Q}))] = \frac{1}{M}\sum_{m=1}^M \phi(\hat{\Delta}_S^n(\mathcal{Q}_m), \hat{\Delta}_Y^n(\mathcal{Q}_m)),
\end{equation}
and the target is:
\begin{equation}
\Theta(\hat{\mu}_M) = \mathbb{E}_{\hat{\mu}_M}[\phi(\Delta_S(\mathcal{Q}), \Delta_Y(\mathcal{Q}))].
\end{equation}

\paragraph{Step 4.2: Decompose estimator (conditional).}
\begin{equation}
\sqrt{n}(\hat{\Theta}_n - \Theta(\hat{\mu}_M)) = \sqrt{n}\left(\mathbb{E}_{\hat{\mu}_M}[\phi(\hat{\Delta}_S^n, \hat{\Delta}_Y^n)] - \mathbb{E}_{\hat{\mu}_M}[\phi(\Delta_S, \Delta_Y)]\right).
\end{equation}

\paragraph{Step 4.3: Expand using functional delta method.}
By Lemma~\ref{lem:functional-expansion}:
\begin{align}
\sqrt{n}(\hat{\Theta}_n - \Theta(\hat{\mu}_M)) &= \sqrt{n} \mathbb{E}_{\hat{\mu}_M}\left[\nabla\phi(\Delta_S, \Delta_Y) \cdot \begin{pmatrix} \hat{\Delta}_S^n - \Delta_S \\ \hat{\Delta}_Y^n - \Delta_Y \end{pmatrix}\right] + o_P(1) \\
&= \mathbb{E}_{\hat{\mu}_M}\left[\nabla\phi(\Delta_S, \Delta_Y) \cdot \begin{pmatrix} \frac{1}{\sqrt{n}}\sum_{i=1}^n \psi_S^{\text{AIPW}}(O_i; \eta_0, \mathcal{Q}) \\ \frac{1}{\sqrt{n}}\sum_{i=1}^n \psi_Y^{\text{AIPW}}(O_i; \eta_0, \mathcal{Q}) \end{pmatrix}\right] + o_P(1).
\end{align}

\paragraph{Step 4.4: Interchange expectation and sum.}
Since $\hat{\mu}_M$ is treated as fixed (conditional on MCMC samples), and by Fubini's theorem:
\begin{align}
\sqrt{n}(\hat{\Theta}_n - \Theta(\hat{\mu}_M)) &= \frac{1}{\sqrt{n}}\sum_{i=1}^n \mathbb{E}_{\hat{\mu}_M}\left[\nabla\phi(\Delta_S, \Delta_Y) \cdot \begin{pmatrix} \psi_S^{\text{AIPW}}(O_i; \eta_0, \mathcal{Q}) \\ \psi_Y^{\text{AIPW}}(O_i; \eta_0, \mathcal{Q}) \end{pmatrix}\right] + o_P(1) \\
&= \frac{1}{\sqrt{n}}\sum_{i=1}^n \psi_\Theta(O_i; \hat{\mu}_M) + o_P(1),
\end{align}
where the composite influence function is:
\begin{equation}
\psi_\Theta(O; \hat{\mu}_M) = \mathbb{E}_{\hat{\mu}_M}\left[\nabla\phi(\Delta_S(\mathcal{Q}), \Delta_Y(\mathcal{Q})) \cdot \begin{pmatrix} \psi_S^{\text{AIPW}}(O; \eta_0, \mathcal{Q}) \\ \psi_Y^{\text{AIPW}}(O; \eta_0, \mathcal{Q}) \end{pmatrix}\right].
\end{equation}

\paragraph{Step 4.5: Apply CLT (conditional).}
Conditional on $\{\mathcal{Q}_1, \ldots, \mathcal{Q}_M\}$, the observations $O_1, \ldots, O_n \overset{\text{i.i.d.}}{\sim} \mathbb{P}_0$ are independent of the MCMC samples. By Assumption~\ref{assm:bounded-if-variance}, $\mathbb{E}_{\mathbb{P}_0}[\psi_\Theta(O; \hat{\mu}_M)^2] < \infty$. By the CLT:
\begin{equation}
\frac{1}{\sqrt{n}}\sum_{i=1}^n \psi_\Theta(O_i; \hat{\mu}_M) \xrightarrow{d} N(0, \mathbb{E}_{\mathbb{P}_0}[\psi_\Theta(O; \hat{\mu}_M)^2]) = N(0, \sigma^2(\hat{\mu}_M)).
\end{equation}

Note that $\mathbb{E}_{\mathbb{P}_0}[\psi_\Theta(O; \hat{\mu}_M)] = 0$ by construction (using Neyman orthogonality from Lemma~\ref{lem:orthogonality}), conditional on the MCMC samples.

\paragraph{Step 4.6: Unconditional result (when $M \to \infty$).}
When $M \to \infty$, by Assumption~\ref{assm:mcmc}, $\hat{\mu}_M \to \mu$ weakly. By Assumption~\ref{assm:bounded-if-variance} (uniform moment bound), $\sup_{\mathcal{Q}} \mathbb{E}[\psi^{\text{AIPW}}(O; \mathcal{Q})^2] < \infty$, hence $\psi_\Theta(O; \hat{\mu}_M) \to \psi_\Theta(O; \mu)$ in $L^2(\mathbb{P}_0)$ by dominated convergence (using bounded $\nabla\phi$ from Assumption~\ref{assm:functional}). Therefore:
\begin{equation}
\sigma^2(\hat{\mu}_M) \to \sigma^2(\mu),
\end{equation}
and unconditionally:
\begin{equation}
\sqrt{n}(\hat{\Theta}_n - \Theta(\mu)) \xrightarrow{d} N(0, \sigma^2(\mu)).
\end{equation}

\paragraph{Conclusion.}
We have shown:
\begin{enumerate}[(a)]
\item Conditional on MCMC samples: $\sqrt{n}(\hat{\Theta}_n - \Theta(\hat{\mu}_M)) \xrightarrow{d} N(0, \sigma^2(\hat{\mu}_M))$.
\item When $M \to \infty$: $\Theta(\hat{\mu}_M) \to \Theta(\mu)$ and $\sigma^2(\hat{\mu}_M) \to \sigma^2(\mu)$.
\end{enumerate}
\end{proof}

\subsection{Variance Estimation}

The asymptotic variance $\sigma^2(\hat{\mu}_M)$ can be consistently estimated (conditional on MCMC samples) by:
\begin{equation}
\hat{\sigma}^2_n = \frac{1}{n}\sum_{i=1}^n \hat{\psi}_\Theta(O_i; \hat{\mu}_M)^2,
\end{equation}
where:
\begin{equation}
\hat{\psi}_\Theta(O_i; \hat{\mu}_M) = \frac{1}{M}\sum_{m=1}^M \nabla\phi(\hat{\Delta}_S^n(\mathcal{Q}_m), \hat{\Delta}_Y^n(\mathcal{Q}_m)) \cdot \begin{pmatrix} \psi_S^{\text{AIPW}}(O_i; \hat{\eta}^{(-k_i)}, \mathcal{Q}_m) \\ \psi_Y^{\text{AIPW}}(O_i; \hat{\eta}^{(-k_i)}, \mathcal{Q}_m) \end{pmatrix}.
\end{equation}
Here, $\hat{\eta}^{(-k_i)}$ denotes the nuisance parameters estimated on the cross-fitting fold that excludes observation $i$. This ensures the influence function estimates are independent of the data used to estimate nuisances, maintaining the double robustness property.

\begin{proposition}[Consistent variance estimation]
\label{prop:variance-consistency}
Under the conditions of Theorem~\ref{thm:asymptotic-normality}, conditional on $\{\mathcal{Q}_1, \ldots, \mathcal{Q}_M\}$:
\begin{equation}
\hat{\sigma}^2_n \xrightarrow{P} \sigma^2(\hat{\mu}_M) \quad \text{as } n \to \infty.
\end{equation}
Furthermore, if $M \to \infty$, then unconditionally $\hat{\sigma}^2_n \xrightarrow{P} \sigma^2(\mu)$.
\end{proposition}

\begin{proof}[Proof of Proposition~\ref{prop:variance-consistency}]
Consistency follows in three steps, plus the unconditional extension.

\paragraph{Step 1: $L^2$ convergence of AIPW influence functions.}
By Assumption~\ref{assm:nuisance} (nuisance convergence rates) and the Lipschitz structure of the AIPW score:
\begin{equation}
\|\psi_S^{\text{AIPW}}(O; \hat{\eta}, \mathcal{Q}) - \psi_S^{\text{AIPW}}(O; \eta_0, \mathcal{Q})\|_{L^2(\mathbb{P}_0)} \xrightarrow{P} 0.
\end{equation}
This follows from:
\begin{itemize}
\item $\|\hat{e} - e\|_{L^2} = o_P(1)$ and $\|\hat{\mu}_a - \mu_a\|_{L^2} = o_P(1)$ (Assumption~\ref{assm:nuisance}(b-c))
\item AIPW score is Lipschitz in $(e, \mu_0, \mu_1)$ on the domain bounded away from $0$ and $1$ (by positivity, Assumption~\ref{assm:nuisance}(a))
\item Bounded weights: $w(X; \mathcal{Q}) \leq C$ (Assumption~\ref{assm:treatment-effects}(c))
\end{itemize}

\paragraph{Step 2: Gradient convergence.}
By Lemma~\ref{lem:cross-fitting}, $\hat{\Delta}_S(\mathcal{Q}_m) \xrightarrow{P} \Delta_S(\mathcal{Q}_m)$ and $\hat{\Delta}_Y(\mathcal{Q}_m) \xrightarrow{P} \Delta_Y(\mathcal{Q}_m)$ for each $m$. By Assumption~\ref{assm:lipschitz} (Lipschitz gradient):
\begin{equation}
\nabla\phi(\hat{\Delta}_S(\mathcal{Q}_m), \hat{\Delta}_Y(\mathcal{Q}_m)) \xrightarrow{P} \nabla\phi(\Delta_S(\mathcal{Q}_m), \Delta_Y(\mathcal{Q}_m)).
\end{equation}

\paragraph{Step 3: Combine via continuous mapping.}
Define $g: L^2(\mathbb{P}_0) \times \mathbb{R}^2 \to \mathbb{R}$ by:
\begin{equation}
g(\psi, \nabla) = \mathbb{E}_{\mathbb{P}_0}[(\nabla \cdot \psi)^2].
\end{equation}
By Steps 1-2, $(\hat{\psi}^{\text{AIPW}}, \nabla\phi(\hat{\Delta})) \xrightarrow{P} (\psi^{\text{AIPW}}, \nabla\phi(\Delta))$ in the product topology. Since $g$ is continuous (by dominated convergence using bounded moments), the continuous mapping theorem gives:
\begin{equation}
\hat{\sigma}^2_n = \frac{1}{n}\sum_{i=1}^n \hat{\psi}_\Theta(O_i)^2 \xrightarrow{P} \mathbb{E}_{\mathbb{P}_0}[\psi_\Theta(O)^2] = \sigma^2(\hat{\mu}_M).
\end{equation}

\paragraph{Step 4: Unconditional convergence when $M \to \infty$.}
When $M \to \infty$, by Assumption~\ref{assm:mcmc}, $\hat{\mu}_M \to \mu$ weakly. By Assumption~\ref{assm:bounded-if-variance} (uniform moment bound), dominated convergence applies, giving $\sigma^2(\hat{\mu}_M) \to \sigma^2(\mu)$. Combining with Step 3 yields unconditional consistency.
\end{proof}

\begin{remark}[Practical variance estimation]
In practice, one estimates $\sigma^2(\hat{\mu}_M)$, which accounts for sampling variation but not MCMC approximation error. If $M$ is large (e.g., $M \geq 500$), the MCMC error is typically negligible relative to sampling uncertainty. Formal accounting for MCMC error would require bootstrap or other resampling methods applied to the MCMC procedure itself, but this is rarely done in practice.
\end{remark}

\subsection{Application to Correlation Functional}\label{sec:correlation-application}

For the correlation functional $\phi(\Delta_S, \Delta_Y) = \text{cor}_\mu(\Delta_S, \Delta_Y)$, this is a composite functional:
\begin{equation}
\phi(\Delta_S, \Delta_Y) = \frac{\mathbb{E}_\mu[\Delta_S \Delta_Y] - \mathbb{E}_\mu[\Delta_S]\mathbb{E}_\mu[\Delta_Y]}{\sqrt{\text{Var}_\mu(\Delta_S) \text{Var}_\mu(\Delta_Y)}}.
\end{equation}

Correlation is thus a smooth function $\varphi$ of the \emph{five base moments} $m(\mu) = (\mathbb{E}_\mu[\Delta_S], \mathbb{E}_\mu[\Delta_Y], \mathbb{E}_\mu[\Delta_S^2], \mathbb{E}_\mu[\Delta_Y^2], \mathbb{E}_\mu[\Delta_S \Delta_Y])$, i.e. $\Theta(\mu) = \varphi(m(\mu))$. This is the object to which the functional delta method applies: $\phi$ is not a fixed map $\mathbb{R}^2 \to \mathbb{R}$ of a single pair $(\Delta_S, \Delta_Y)$, but a functional of the across-study \emph{distribution} $\mu$ of the pair, expressed through these five $\mu$-moments. Each base moment is a smooth (polynomial) map of $(\Delta_S(\mathcal{Q}), \Delta_Y(\mathcal{Q}))$ averaged over $\mathcal{Q} \sim \mu$, so the per-study gradient $\nabla\phi(\Delta_S(\mathcal{Q}_m), \Delta_Y(\mathcal{Q}_m))$ used in \eqref{eq:aipw-s}--\eqref{eq:aipw-y} and Lemma~\ref{lem:functional-expansion} is the chain-rule composition of $\nabla\varphi(m(\mu))$ with the moment derivatives.

\begin{remark}
The map $\varphi$ is continuously differentiable, hence Hadamard differentiable, on the region $\{\mathrm{Var}_\mu(\Delta_S) \geq c_0,\ \mathrm{Var}_\mu(\Delta_Y) \geq c_0\}$ guaranteed by Assumption~\ref{assm:nondegeneracy}; its gradient is bounded there (scaling like $c_0^{-3/2}$). This makes the non-degeneracy requirement explicit rather than implicit in a ``bounded gradient'' assumption, and localizes where the delta-method linearization degrades (as $c_0 \to 0$).
\end{remark}

\subsection{Post-Proof Audit}

\paragraph{Assumption check.}
All assumptions stated upfront in Section 2.2 and explicitly invoked throughout:
\begin{itemize}
\item Assumptions~\ref{assm:treatment-effects}--\ref{assm:functional}: Treatment effects and functional regularity
\item Assumptions~\ref{assm:nuisance}--\ref{assm:lipschitz}: Nuisance estimation (new for AIPW)
\item Assumptions~\ref{assm:mcmc}--\ref{assm:geometry}: MCMC and geometry
\item Assumption~\ref{assm:sample-size}: Sample size
\end{itemize}

\paragraph{Dependency check.}
All invoked results have their conditions verified:
\begin{itemize}
\item \textbf{Chernozhukov et al.\ (2018) Theorem 3.1:} Conditions verified in Lemma~\ref{lem:cross-fitting} (Neyman orthogonality, $L^2$ rates, cross-fitting, bounded moments)
\item \textbf{CLT:} i.i.d.\ observations, finite variance (Assumption~\ref{assm:bounded-if-variance})
\item \textbf{Delta method:} Hadamard differentiability (Assumption~\ref{assm:functional}), Lipschitz gradient (Assumption~\ref{assm:lipschitz})
\item \textbf{Fubini:} Bounded integrands (Assumptions~\ref{assm:treatment-effects}(c), \ref{assm:functional}, \ref{assm:bounded-if-variance})
\item \textbf{Dominated convergence:} Uniform moment bounds (Assumption~\ref{assm:bounded-if-variance})
\end{itemize}

\paragraph{Rate verification.}
\begin{itemize}
\item \textbf{Nuisance estimation:} $o_P(n^{-1/4})$ for propensity and outcome regressions (Assumption~\ref{assm:nuisance}(b-c))
\item \textbf{Second-order remainder:} $\|\hat{\eta} - \eta_0\|_{L^2}^2 = [o_P(n^{-1/4})]^2 = o_P(n^{-1/2})$ (Assumption~\ref{assm:nuisance}(d))
\item \textbf{Treatment effect estimators:} $O_P(n^{-1/2})$ rate (Lemma~\ref{lem:cross-fitting})
\item \textbf{MCMC approximation:} $O_P(M^{-1/2})$, independent of $n$ (Assumption~\ref{assm:mcmc})
  \begin{itemize}
  \item If $M$ fixed: Non-vanishing bias in target
  \item If $M \to \infty$: Bias vanishes
  \end{itemize}
\item \textbf{Remainder terms:} All $o_P(n^{-1/2})$ by Neyman orthogonality and bounded second derivatives
\item \textbf{Final rate:} $\sqrt{n}$-consistent (conditional on MCMC samples)
\end{itemize}

\paragraph{Key differences from RCT case.}
\begin{itemize}
\item \textbf{Influence function:} Three-term AIPW structure replaces simple weighted difference-in-means
\item \textbf{Nuisance estimation:} Propensity scores and outcome regressions estimated with $o_P(n^{-1/4})$ rates
\item \textbf{Cross-fitting:} Sample splitting prevents overfitting bias
\item \textbf{Neyman orthogonality:} First-order nuisance bias eliminated (Lemma~\ref{lem:orthogonality})
\item \textbf{Double robustness:} Consistency if either $\hat{e}$ or $\hat{\mu}_a$ correct (structural property of AIPW)
\end{itemize}

\paragraph{What stays the same.}
\begin{itemize}
\item MCMC approximation rate: Still $O_P(M^{-1/2})$
\item Functional delta method: Still applies via Hadamard differentiability
\item Final $\sqrt{n}$ rate: Unchanged
\item Bounded weights key to uniformity: Still central to proof
\end{itemize}

\paragraph{Edge case remarks.}
\begin{enumerate}
\item \textbf{Degenerate variance (correlation functional):} If $\text{Var}_\mu(\Delta_S) = 0$ or $\text{Var}_\mu(\Delta_Y) = 0$, the correlation functional is undefined. The proof requires non-degeneracy: $\min(\text{Var}_\mu(\Delta_S), \text{Var}_\mu(\Delta_Y)) \geq c_0 > 0$, stated explicitly as Assumption~\ref{assm:nondegeneracy}.

\item \textbf{Extreme reweighting:} If $\max_i w_i(\mathcal{Q}) \approx C$ (near the bound), variance can inflate significantly. In practice, recommend $C \leq 10$ to avoid extreme variance inflation. Effective sample size drops as $n_{\text{eff}} \approx n / (1 + \text{CV}^2)$ where $\text{CV}$ is coefficient of variation of weights.

\item \textbf{Non-smooth functionals:} Indicator functions (e.g., $\phi(\Delta_S, \Delta_Y) = \mathbb{1}(\Delta_S > 0)$) are not Hadamard differentiable. The proof applies only to smooth functionals: correlation, covariance, products of moments. For non-smooth functionals, bootstrap or other methods required.

\item \textbf{$M \to \infty$ clarification:} Three separate issues with large $M$:
\begin{itemize}
\item \textbf{Remainder uniformity:} No restriction on $M$ growth rate (bound independent of $M$). This is the key innovation.
\item \textbf{CLT for composite IF:} Requires uniform moment bound $\sup_{\mathcal{Q}} \mathbb{E}[\psi^2] < \infty$ (Assumption~\ref{assm:bounded-if-variance}). Without this, CLT can break as $M \to \infty$.
\item \textbf{MCMC approximation quality:} Requires $\|\hat{\mu}_M - \mu\|_{\text{TV}} = O_P(M^{-1/2})$ (Assumption~\ref{assm:mcmc}). About approximation error, not uniformity.
\end{itemize}

\item \textbf{Strong positivity threshold:} Assumption~\ref{assm:nuisance}(a) requires $e(X) \in [0.1, 0.9]$ (stronger than standard $e(X) \in (0, 1)$). This ensures second derivative bound $\frac{2M}{c^3} \leq 2000$ stays reasonable. Near boundaries ($e(X) \approx 0$ or $1$), second derivatives explode, violating Assumption~\ref{assm:second-derivatives}. Use trimming or alternative methods for weak overlap regions.
\end{enumerate}

\paragraph{Tightness assessment.}
The $\sqrt{n}$ rate is sharp (semiparametric efficiency bound for this estimand). The $o_P(n^{-1/4})$ nuisance rates are also sharp under standard smoothness conditions (cannot be weakened without additional structure). Assumptions are near-minimal for these rates.

\paragraph{Failure localization.}
If asymptotic normality fails or coverage is poor:
\begin{enumerate}
\item \textbf{Check propensity overlap:} $\min_i \hat{e}(X_i) \geq 0.1$ and $\max_i \hat{e}(X_i) \leq 0.9$? If not, trim or use alternative methods.
\item \textbf{Check nuisance convergence:} Use cross-validated MSE to verify $o_P(n^{-1/4})$ rates. If slower, coverage will be poor.
\item \textbf{Check weight bounds:} $\max_m \max_i w_i(\mathcal{Q}_m) \leq C$? Large weights inflate variance.
\item \textbf{Check variance degeneracy:} $\min(\text{Var}_\mu(\Delta_S), \text{Var}_\mu(\Delta_Y)) > 0$? If near-zero, correlation undefined.
\item \textbf{Check MCMC convergence:} Is $M$ large enough? Is burn-in sufficient? Use $\hat{R}$ diagnostics.
\item \textbf{Check MCMC bias:} If $M$ is fixed and coverage systematically off, MCMC bias $O_P(M^{-1/2})$ may dominate. Increase $M$.
\end{enumerate}

\paragraph{Implementation notes.}
The R package \texttt{surrogateTransportability} implements this framework:
\begin{itemize}
\item \texttt{tv\_ball\_correlation\_IF\_adaptive()}: main estimator (adaptive hit-and-run,
  importance-weighting or cross-fitted AIPW, influence-function inference)
\item \texttt{sample\_tv\_ball()}: uniform hit-and-run sampling on the TV ball
\item \texttt{gradient\_correlation\_analytical()}: analytic gradient of the correlation functional
\end{itemize}
See package documentation for usage examples.

% =====================================================================
%  GENERAL ASYMPTOTIC-NORMALITY THEORY (extends Theorem \ref{thm:asymptotic-normality})
%  Labels use an "A" scheme (thm:general, lem:*-A, cor:finite-support).
%  Notation matches the finite-support proof above (explicit long form).
% =====================================================================

\section{General Estimand: Asymptotic Normality via the Bilinear-Functional Structure}
\label{sec:general-theory}

This section extends Theorem~\ref{thm:asymptotic-normality} from the finite-support,
correlation-specific case to a general smooth functional $\Theta = \Psi(m_1,\dots,m_K)$
of the across-study effect distribution, on discrete \emph{or} smooth-enough continuous
covariate spaces. As in Theorem~\ref{thm:asymptotic-normality}, we \emph{condition on the
sampled geometry}: the MCMC sample $\hat{\mu}_M$ (equivalently the reweighting-covariance
kernel $C$ it induces, defined below) is treated as fixed, and the $\sqrt{M}$ MCMC term is
handled exactly as before via Lemma~\ref{lem:mcmc-error}.

\subsection{General estimand and the bilinear-functional structure}
\label{subsec:bilinear-structure}

\paragraph{Covariate-level effects.}
Let $O=(X,A,S,Y)\sim\mathbb{P}_0$ with $A\in\{0,1\}$ and $X\in\mathcal{X}$, where
$\mathcal{X}\subseteq\mathbb{R}^d$ (finite, or continuous of dimension $d$). Define the
conditional average treatment effects (CATEs)
\begin{equation}
\tau_S(x) = \mathbb{E}_{\mathbb{P}_0}[S(1)-S(0)\mid X=x],
\qquad
\tau_Y(x) = \mathbb{E}_{\mathbb{P}_0}[Y(1)-Y(0)\mid X=x],
\end{equation}
and write $\beta=(\tau_S,\tau_Y)$. For a future study $\mathcal{Q}\ll\mathbb{P}_0$ with
density ratio $w(x)=\tfrac{d\mathcal{Q}}{d\mathbb{P}_0}(x)$, the study-level effects are the
\emph{linear} functionals
\begin{equation}
\Delta_a(\mathcal{Q}) = \mathbb{E}_{\mathbb{P}_0}[\,w(X)\,\tau_a(X)\,],
\qquad a\in\{S,Y\}.
\label{eq:delta-linear}
\end{equation}

\paragraph{The geometry and its kernel.}
A geometry is a probability measure $\mu$ over reweightings $w$, supported on a local
convex body $\{\mathcal{Q}:d(\mathcal{Q},\mathbb{P}_0)\le\lambda\}$ (Assumption~\ref{assm:geometry}).
Define
\begin{equation}
\bar{a}(x) = \mathbb{E}_\mu[w(x)],
\qquad
C(x,x') = \mathrm{Cov}_\mu\!\big(w(x),\,w(x')\big),
\label{eq:kernel-def}
\end{equation}
the reweighting mean and the \emph{reweighting-covariance kernel}. By construction $C$ is
symmetric; under Assumptions~\ref{assm:treatment-effects}(c) and \ref{assm:geometry} it is
bounded, $\sup_{x,x'}|C(x,x')|\le C_w^2<\infty$.

\paragraph{Linear/bilinear reduction.}
Taking $\mu$-moments of \eqref{eq:delta-linear} and using Fubini (justified by bounded weights
and $\mathbb{E}_{\mathbb{P}_0}[\tau_a^2]<\infty$):
\begin{align}
\mathbb{E}_\mu[\Delta_a(\mathcal{Q})]
   &= \int \bar{a}(x)\,\tau_a(x)\,d\mathbb{P}_0(x),
   \tag{linear}\label{eq:reduction-linear}\\
\mathrm{Cov}_\mu\!\big(\Delta_a(\mathcal{Q}),\Delta_b(\mathcal{Q})\big)
   &= \iint C(x,x')\,\tau_a(x)\,\tau_b(x')\,d\mathbb{P}_0(x)\,d\mathbb{P}_0(x')
   \;=:\; \Theta_{ab},
   \tag{bilinear}\label{eq:reduction-bilinear}
\end{align}
for $a,b\in\{S,Y\}$. Thus every load-bearing $\mu$-moment is either a bounded \emph{linear}
functional of $\beta$ (against the weight $\bar{a}$) or a bounded \emph{bilinear} functional
of $\beta$ (against the kernel $C$). We write the moment vector
\begin{equation}
m = \big(\underbrace{\mathbb{E}_\mu\Delta_S,\ \mathbb{E}_\mu\Delta_Y}_{\text{linear}},\ 
        \underbrace{\Theta_{SS},\ \Theta_{YY},\ \Theta_{SY}}_{\text{bilinear}},\ \dots\big)
        \in\mathbb{R}^K,
\end{equation}
and the estimand is $\Theta = \Psi(m)$ for a smooth $\Psi:\mathbb{R}^K\to\mathbb{R}$
(Assumption~\ref{assm:functional}). Correlation, $R^2$, MSPE, and concordance are instances.

\paragraph{Discrete specialization ($C\to\Sigma$).}
When $X\in\{1,\dots,K\}$ is finite with cell probabilities $p_0(k)$, write
$\tau_a=(\tau_a(1),\dots,\tau_a(K))^\top$ and $q=(q_1,\dots,q_K)^\top$ the future-study cell
mass. Then $w(k)=q_k/p_0(k)$, $\bar{a}(k)=\mathbb{E}_\mu[w(k)]$, and the kernel collapses to a
$K\times K$ matrix $\Sigma$ with entries
$\Sigma_{kk'} = C(k,k')\,p_0(k)\,p_0(k')=\mathrm{Cov}_\mu(q_k,q_{k'})$, so that
\begin{equation}
\Theta_{ab} = \iint C\,\tau_a\tau_b\,d\mathbb{P}_0^2
            = \sum_{k,k'} p_0(k)p_0(k')\,C(k,k')\,\tau_a(k)\tau_b(k')
            = \tau_a^\top \Sigma\, \tau_b .
\label{eq:discrete-quadform}
\end{equation}
This is the $\tau^\top\Sigma\tau$ object of the finite-support proof; Corollary~\ref{cor:finite-support}
shows the general result collapses to Theorem~\ref{thm:asymptotic-normality} here.

\subsection{Assumption ledger A0--A8}
\label{subsec:assumptions-A}

Assumptions \ref{assm:treatment-effects}--\ref{assm:functional} and
\ref{assm:nuisance}--\ref{assm:sample-size} carry over as the general-case A1--A4, A6--A8
(with $\phi$-specific statements reread for $\Psi$ and the moment vector $m$). We state in full
the two assumptions that are new relative to the finite-support proof: A0 (the causal scope of
the estimand) and A5 (the bilinear-functional estimability boundary).

\begin{assumption}[A0: Conditional-effect transportability given $X$]
\label{assm:A0-transport}
All cross-study variation in treatment effects is compositional in the measured covariates
$X$: future studies $\mathcal{Q}$ differ from $\mathbb{P}_0$ only in the marginal law of $X$
and share $\mathbb{P}_0$'s conditional average treatment effects $\tau_S(\cdot),\tau_Y(\cdot)$
given $X$. Equivalently, the reweighting $w=\tfrac{d\mathcal{Q}}{d\mathbb{P}_0}$ is a function
of $X$ alone, so $\tau_S(x),\tau_Y(x)$ are frozen at their $\mathbb{P}_0$ values under every
$\mathcal{Q}$, which is exactly what makes \eqref{eq:delta-linear} hold. (Only the CATEs, not
the full conditional law of potential outcomes, need be shared: \eqref{eq:delta-linear} depends
on $\mathbb{P}_0$ through $(\tau_S,\tau_Y)$ alone.)
\end{assumption}

\begin{remark}[Definitional vs.\ interpretive reading of A0]
\label{rem:A0-reading}
Under the \emph{definitional} reading the estimand simply \emph{is} the functional over the
class of $X$-compositional future studies, and A0 is a scope statement, not an assumption;
under the \emph{interpretive} reading---reading $\Theta$ as ``does the surrogate transport to
real future studies''---A0 is a substantive causal assumption of the effects-transport-given-$X$
class, distinct from the support condition $\mathcal{Q}\ll\mathbb{P}_0$
(Assumption~\ref{assm:treatment-effects}(c)) and not removable for a single-study point estimate.
\end{remark}

\begin{assumption}[A5: Bilinear-functional $\sqrt{n}$-estimability]
\label{assm:A5-elbow}
Each bilinear functional $\Theta_{ab}$ in \eqref{eq:reduction-bilinear} admits a
$\sqrt{n}$-consistent, asymptotically linear (debiased) estimator. Let $\tau_S,\tau_Y$ lie in
H\"older/Sobolev balls of smoothness $s_S,s_Y$ on $\mathcal{X}\subseteq\mathbb{R}^d$. Then, for
the identity/Dirac kernel, $\sqrt{n}$-estimability of $\Theta_{ab}$ holds if and only if the
sum-of-smoothness condition
\begin{equation}
s_S + s_Y > d/2
\label{eq:elbow}
\end{equation}
holds. When \eqref{eq:elbow} fails, the minimax rate is
$n^{-\,4s/(4s+d)}$ (with $s=\tfrac{1}{2}(s_S+s_Y)$ the effective smoothness; this single-$s$
rate expression is the balanced-smoothness reduction of the genuinely two-smoothness minimax
boundary, whereas the elbow \eqref{eq:elbow} itself is exact), strictly slower
than $n^{-1/2}$, and no $\sqrt{n}$-consistent asymptotically normal estimator exists
\cite{BickelRitov1988,BirgeMassart1995,Robins2008HOIF,Robins2017Minimax,Kennedy2024Minimax,McClean2024DCDR}.
\end{assumption}

\begin{remark}[A5 is conservative; finite support is automatic]
\label{rem:A5-conservative}
Condition \eqref{eq:elbow} is the boundary for the \emph{Dirac} kernel. Our kernel $C$ is
bounded and (typically) smooth under Assumptions~\ref{assm:treatment-effects}(c) and
\ref{assm:geometry}; a genuinely smooth $C$ typically \emph{absorbs} smoothness (its integral
operator regularizes $h_b=\int C\tau_b\,d\mathbb{P}_0$) and thus \emph{relaxes} the requirement,
so \eqref{eq:elbow} is, at worst, a conservative sufficient condition here. On finite
$X$ every function is arbitrarily smooth and $\Theta_{ab}=\tau_a^\top\Sigma\tau_b$ is a fixed
smooth (indeed polynomial) function of finitely many cell means, so A5 holds automatically---%
this is why the finite-support proof never states it. Higher-order influence functions attain
the minimax rate below the elbow but do \emph{not} restore $\sqrt{n}$ normality.
\end{remark}

\subsection{Layer L1: CATE asymptotic linearity}

\begin{lemma}[A1: CATE asymptotic linearity]
\label{lem:cate-linearity-A}
Under A1--A2 (Assumption~\ref{assm:treatment-effects} generalized off the finite simplex, plus
a pluggable estimator supplying a CATE influence function), for each $a\in\{S,Y\}$ and each
$x\in\mathcal{X}$ the CATE estimator admits the expansion
\begin{equation}
\hat{\tau}_a(x) - \tau_a(x)
  = \frac{1}{n}\sum_{i=1}^n \mathrm{IF}_{\tau_a}(O_i;x) + R_a(x),
\label{eq:cate-expansion}
\end{equation}
where $\mathbb{E}_{\mathbb{P}_0}[\mathrm{IF}_{\tau_a}(O;x)]=0$ for every $x$, and the remainder
satisfies $\|R_a\|_{L^2(\mathbb{P}_0)}=o_P(n^{-1/4})$ (product-rate control under A6).
For an RCT $\mathrm{IF}_{\tau_a}$ is the (per-$x$) difference-of-means score; for AIPW /
DR-learners it is the pointwise AIPW score, and a generic learner is admissible whenever
it supplies a valid $\mathrm{IF}_{\tau_a}$ satisfying \eqref{eq:cate-expansion}.
\end{lemma}

\begin{proof}[Proof of Lemma~\ref{lem:cate-linearity-A}]
This is the pointwise (per-$x$) restatement of the AIPW expansion established, integrated
against $w$, in Lemma~\ref{lem:cross-fitting}; the reweighted per-study score
$\psi_a^{\text{AIPW}}(O;\eta_0,\mathcal{Q})$ of \eqref{eq:aipw-s}--\eqref{eq:aipw-y} equals
$\int w(x)\,\mathrm{IF}_{\tau_a}(O;x)\,d\mathbb{P}_0(x)$ up to the mean-zero centering, so
mean-zero-ness and the $o_P(n^{-1/2})$ integrated remainder transfer directly. For the discrete
case the score is exactly the per-cell AIPW score of the companion derivation
(\texttt{derivation\_influence\_functions.md}, \S2). The remainder rate is A6.
% NOTE(auditor): \eqref{eq:cate-expansion} is stated pointwise in x. Downstream (Lemma
% \ref{lem:bilinear-eif-A}) we only ever integrate IF against dP0 or the bounded kernel C, so a
% pointwise-in-x, L2-in-O influence function suffices; no Donsker/uniform-in-x class is claimed.
\end{proof}

\subsection{Layer L3: the bilinear-functional EIF (core new result)}

\begin{lemma}[A3: EIF of a bilinear functional of CATEs]
\label{lem:bilinear-eif-A}
Fix a bounded symmetric kernel $C$ and the covariate law of $\mathbb{P}_0$; treat both as
known (conditioning on the sampled geometry, as in Theorem~\ref{thm:asymptotic-normality}).
Suppose $\tau_a,\tau_b$ are pathwise differentiable with influence functions
$\mathrm{IF}_{\tau_a},\mathrm{IF}_{\tau_b}$ (Lemma~\ref{lem:cate-linearity-A}), and
$\mathbb{E}_{\mathbb{P}_0}[\tau_a^2],\mathbb{E}_{\mathbb{P}_0}[\tau_b^2]<\infty$. Then
$\Theta_{ab}=\iint C(x,x')\tau_a(x)\tau_b(x')\,d\mathbb{P}_0(x)d\mathbb{P}_0(x')$ is pathwise
differentiable with (mean-zero) efficient influence function
\begin{equation}
\chi_{ab}(O)
 = \int h_b(x)\,\mathrm{IF}_{\tau_a}(O;x)\,d\mathbb{P}_0(x)
 + \int h_a(x')\,\mathrm{IF}_{\tau_b}(O;x')\,d\mathbb{P}_0(x'),
\label{eq:chi-ab}
\end{equation}
where
\begin{equation}
h_b(x) = \int C(x,x')\,\tau_b(x')\,d\mathbb{P}_0(x'),
\qquad
h_a(x') = \int C(x',x)\,\tau_a(x)\,d\mathbb{P}_0(x).
\label{eq:h-directions}
\end{equation}
Since $\mathbb{E}_{\mathbb{P}_0}[\mathrm{IF}_{\tau_a}(O;x)]=0$ for every $x$,
$\mathbb{E}_{\mathbb{P}_0}[\chi_{ab}]=0$; the plug-in value $\Theta_{ab}$ is carried by the
plug-in term of the one-step estimator \eqref{eq:onestep}, not by the influence function.
In particular, when $a=b$,
$\chi_{aa}(O) = 2\int h_a(x)\,\mathrm{IF}_{\tau_a}(O;x)\,d\mathbb{P}_0(x)$ with
$h_a(x)=\int C(x,x')\tau_a(x')d\mathbb{P}_0(x')$; and for finite $X$, \eqref{eq:chi-ab}
reduces to a quadratic/trace form in $\Sigma$ (Corollary~\ref{cor:finite-support}).
\end{lemma}

\begin{proof}[Proof of Lemma~\ref{lem:bilinear-eif-A}]
\emph{Strategy.} $\Theta_{ab}$ depends on $\mathbb{P}_0$ (with $C$ and the covariate law
fixed) only through the pair $(\tau_a,\tau_b)$. We compute the pathwise derivative by the
product rule over the two $\tau$-factors, identify the Riesz representer of each directional
derivative, and symmetrize.

\emph{Step 1: one-dimensional submodel.} Let $\{\mathbb{P}_t\}$ be a regular parametric submodel
through $\mathbb{P}_0$ at $t=0$ with score $g(O)$, and let $\tau_{a,t},\tau_{b,t}$ denote the
CATEs under $\mathbb{P}_t$. Pathwise differentiability of the CATEs (Lemma~\ref{lem:cate-linearity-A})
gives, for each fixed $x$,
\begin{equation}
\frac{d}{dt}\tau_{a,t}(x)\Big|_{t=0}
   = \mathbb{E}_{\mathbb{P}_0}\!\big[\mathrm{IF}_{\tau_a}(O;x)\,g(O)\big],
\label{eq:tau-derivative}
\end{equation}
and similarly for $\tau_b$. (Here $C$ and $d\mathbb{P}_0$ in the definition of $\Theta_{ab}$ are
held fixed; see Remark~\ref{rem:kernel-dependence}.)

\emph{Step 2: product rule over the two factors.} Write
$\Theta_{ab}(t)=\iint C(x,x')\tau_{a,t}(x)\tau_{b,t}(x')\,d\mathbb{P}_0(x)d\mathbb{P}_0(x')$
(covariate law fixed). Differentiating under the integral (Leibniz; justified because $C$ is
bounded, $\tau_{a,t},\tau_{b,t}$ have square-integrable derivatives, and the domination is by
$C_w^2(\,|\partial_t\tau_a||\tau_b| + |\tau_a||\partial_t\tau_b|)\in L^1$):
\begin{align}
\frac{d}{dt}\Theta_{ab}(t)\Big|_{t=0}
 &= \iint C(x,x')\,\Big[\dot\tau_a(x)\,\tau_b(x') + \tau_a(x)\,\dot\tau_b(x')\Big]
     \,d\mathbb{P}_0(x)d\mathbb{P}_0(x') \nonumber\\
 &= \int \dot\tau_a(x)\,h_b(x)\,d\mathbb{P}_0(x)
  + \int \dot\tau_b(x')\,h_a(x')\,d\mathbb{P}_0(x'),
\label{eq:productrule}
\end{align}
where $\dot\tau_a(x)=\tfrac{d}{dt}\tau_{a,t}(x)|_{t=0}$ and we used Fubini (bounded $C$,
square-integrable $\tau$) to collapse each double integral, defining the two \emph{directional}
representers $h_b,h_a$ of \eqref{eq:h-directions}. This is the ``two directional derivatives,
then symmetrize'' step: the $a$-factor is differentiated against the $b$-weight $h_b$, and the
$b$-factor against the $a$-weight $h_a$.

\emph{Step 3: substitute the CATE scores.} Insert \eqref{eq:tau-derivative} into
\eqref{eq:productrule} and interchange the (finite-measure, bounded-integrand) $x$- and
$O$-integrals by Fubini:
\begin{align}
\frac{d}{dt}\Theta_{ab}(t)\Big|_{t=0}
 &= \mathbb{E}_{\mathbb{P}_0}\!\Big[\Big(\int h_b(x)\mathrm{IF}_{\tau_a}(O;x)d\mathbb{P}_0(x)
     + \int h_a(x')\mathrm{IF}_{\tau_b}(O;x')d\mathbb{P}_0(x')\Big)\,g(O)\Big].
\label{eq:riesz}
\end{align}
\emph{Step 4: identify the EIF.} Equation \eqref{eq:riesz} exhibits the pathwise derivative as
$\mathbb{E}_{\mathbb{P}_0}[\chi_{ab}(O)g(O)]$ for every score $g$, where $\chi_{ab}$ is the
bracketed term of \eqref{eq:chi-ab}; by the Riesz representation of pathwise derivatives
$\chi_{ab}$ is the (canonical-gradient) influence function of $\Theta_{ab}$. Since
$\mathbb{E}_{\mathbb{P}_0}[\mathrm{IF}_{\tau_a}(O;x)]=0$ for all $x$, $\chi_{ab}$ is mean-zero,
as an influence function must be; the plug-in value $\Theta_{ab}$ enters not $\chi_{ab}$ but the
plug-in term of the one-step estimator \eqref{eq:onestep}. Efficiency holds because, in the
nonparametric model where the tangent space is all of $L^2_0(\mathbb{P}_0)$, $\chi_{ab}$ is a
$\mathbb{P}_0$-integral of the efficient CATE influence functions and hence lies in the tangent
space, so it is the canonical gradient.

\emph{Step 5: diagonal check.} If $a=b$ then $h_a=h_b$ and $\mathrm{IF}_{\tau_a}=\mathrm{IF}_{\tau_b}$,
so \eqref{eq:chi-ab} becomes
$\chi_{aa}(O)=2\int h_a(x)\mathrm{IF}_{\tau_a}(O;x)d\mathbb{P}_0(x)$, the factor $2$
being exactly the two identical directional derivatives---consistent with $\Theta_{aa}$ being a
pure quadratic whose derivative carries the classical factor $2$.

\emph{Step 6: discrete check.} On finite $X$, \eqref{eq:h-directions} gives
$h_b(k)=\sum_{k'}C(k,k')\tau_b(k')p_0(k')$; substituting into \eqref{eq:chi-ab} and using
$\int h_b\,\mathrm{IF}_{\tau_a}d\mathbb{P}_0=\sum_k p_0(k)h_b(k)\mathrm{IF}_{\tau_a}(\cdot;k)$
reproduces the discrete $\Sigma$-quadratic score; contracting against the CATE-score covariance
$V$ yields the $\mathrm{tr}(\Sigma V)$ debiasing term of the finite-support case
(Corollary~\ref{cor:finite-support}).
\end{proof}

\begin{remark}[Scope: kernel and covariate-law dependence held fixed]
\label{rem:kernel-dependence}
Lemma~\ref{lem:bilinear-eif-A} differentiates only through $(\tau_a,\tau_b)$, holding the
kernel $C$ and the covariate law $d\mathbb{P}_0$ fixed. This matches
Theorem~\ref{thm:asymptotic-normality}, which conditions on the sampled geometry $\hat{\mu}_M$
(hence on $C$). The \emph{unconditional} EIF, which would additionally differentiate $C$'s and
the covariate law's dependence on $\mathbb{P}_0$, contributes further terms; deriving it is
future work and is not needed for the conditional statement of Theorem~\ref{thm:general}.
% NOTE(auditor): Genuine deferral, not a hidden step. Conditional-on-geometry stance inherited
% from Theorem \ref{thm:asymptotic-normality}. An unconditional statement would add Gateaux terms
% from differentiating dP0(x), dP0(x'), and (if C is a P0-plug-in) C itself. Localize here.
\end{remark}

\subsection{Layer L3: one-step debiased cross-fit estimator}

\begin{lemma}[A4: debiased cross-fit estimator of $\Theta_{ab}$]
\label{lem:onestep-A}
Let $\hat\tau_a,\hat\tau_b,\hat{C}$ be cross-fitted estimators and
$\hat h_b(x)=\int \hat C(x,x')\hat\tau_b(x')d\mathbb{P}_0(x')$ (etc.). The one-step estimator
\begin{equation}
\hat\Theta_{ab}
 = \iint \hat C(x,x')\hat\tau_a(x)\hat\tau_b(x')\,d\mathbb{P}_0(x)d\mathbb{P}_0(x')
 + \frac{1}{n}\sum_{i=1}^n \hat\chi_{ab}(O_i)
\label{eq:onestep}
\end{equation}
is Neyman-orthogonal in $\beta=(\tau_a,\tau_b)$: the Gateaux derivative of the map
$\beta\mapsto \mathbb{E}_{\mathbb{P}_0}[\text{(plug-in)}+\chi_{ab}]$ vanishes at the truth.
\end{lemma}

\begin{proof}[Proof of Lemma~\ref{lem:onestep-A}]
By Lemma~\ref{lem:bilinear-eif-A}, $\chi_{ab}$ is the canonical gradient of $\Theta_{ab}$ in
$\beta$; adding it to the plug-in cancels the first-order plug-in bias, the defining property of
the one-step construction. Concretely, the directional derivative of the plug-in at $\beta$
along $(\delta_a,\delta_b)$ is $\int h_b\delta_a\,d\mathbb{P}_0+\int h_a\delta_b\,d\mathbb{P}_0$
by \eqref{eq:productrule}, while $\mathbb{E}_{\mathbb{P}_0}[\chi_{ab}]$ has directional derivative
equal to its negative; the two cancel. The representer $h_b$ plays the role of the weight $w$ in
Lemma~\ref{lem:orthogonality}: since $h_b,h_a$ are functions of $x$ that do \emph{not} depend on
the outcome-model nuisances $\eta=(e,\mu_0,\mu_1)$ (only on $C$ and $\tau$), the
``orthogonality survives reweighting'' argument of Lemma~\ref{lem:orthogonality}, Step~4, applies
verbatim with $h_b$ (resp.\ $h_a$) in the role of $w(X;\mathcal{Q})$, so orthogonality \emph{in
$\eta$} holds in addition to orthogonality \emph{in $\beta$}. (The estimated $\hat h_b=\int\hat
C\hat\tau_b\,d\mathbb{P}_0$ depends on $\eta$ only through $\hat\tau_b$, which is the
$\beta$-argument handled by $\beta$-orthogonality, not through a separate $\eta$-path, since
$\hat C$ is geometry and hence $\eta$-free; the inheritance is therefore airtight.)
% NOTE(auditor): Two orthogonalities: (i) in beta by the one-step construction; (ii) in eta,
% inherited from Lemma \ref{lem:orthogonality} because h_a,h_b are eta-free x-functions. Both are
% needed for the Lemma \ref{lem:elbow} remainder to be second order. Most worth re-checking.
\end{proof}

\subsection{Layer L3: remainder and the estimability elbow}

\begin{lemma}[A5: bilinear remainder / elbow]
\label{lem:elbow}
Under Lemmas~\ref{lem:cate-linearity-A}, \ref{lem:onestep-A}, the estimator
\eqref{eq:onestep} satisfies
$\hat\Theta_{ab} - \Theta_{ab} = \frac{1}{n}\sum_{i=1}^n \chi_{ab}(O_i) + R_n$,
where the second-order remainder is, up to $o_P(n^{-1/2})$ nuisance terms controlled by A6 and
Lemma~\ref{lem:onestep-A},
\begin{equation}
R_n = \iint \hat C(x,x')\,(\hat\tau_a-\tau_a)(x)\,(\hat\tau_b-\tau_b)(x')\,d\mathbb{P}_0(x)d\mathbb{P}_0(x').
\label{eq:remainder}
\end{equation}
By Cauchy--Schwarz with the bounded symmetric kernel $\hat C$ (operator norm
$\|\hat C\|_{op}\le C_w^2$ under Assumptions~\ref{assm:treatment-effects}(c), \ref{assm:geometry}),
\begin{equation}
|R_n|
 \le \|\hat C\|_{op}\,\big\|\hat\tau_a-\tau_a\big\|_{L^2(\mathbb{P}_0)}\,
     \big\|\hat\tau_b-\tau_b\big\|_{L^2(\mathbb{P}_0)}.
\label{eq:CS-remainder}
\end{equation}
Hence $R_n=o_P(n^{-1/2})$ if $\|\hat\tau_a-\tau_a\|\cdot\|\hat\tau_b-\tau_b\| = o_P(n^{-1/2})$,
which holds under A5: for H\"older/Sobolev smoothness $s_S,s_Y$ the achievable rates give a
product $o_P(n^{-1/2})$ exactly when $s_S+s_Y>d/2$. When \eqref{eq:elbow} fails, $\hat\Theta_{ab}$
has the slower minimax rate $n^{-4s/(4s+d)}$ and no $\sqrt{n}$-normal estimator exists
\cite{BickelRitov1988,BirgeMassart1995,Robins2008HOIF,Kennedy2024Minimax,McClean2024DCDR}.
\end{lemma}

\begin{proof}[Proof of Lemma~\ref{lem:elbow}]
The von Mises expansion of \eqref{eq:onestep} has exact second-order remainder
\eqref{eq:remainder}: the one-step correction removes the first-order term
(Lemma~\ref{lem:onestep-A}), and $\Theta_{ab}$ is exactly \emph{bilinear} in $\beta$, so the
Taylor expansion terminates at the product term with \emph{no} higher-order terms (a quadratic
functional has zero third derivative). The nuisance-$\eta$ contribution is $o_P(n^{-1/2})$ by
Neyman orthogonality in $\eta$ (Lemma~\ref{lem:onestep-A}) and the product-rate A6, exactly as in
Lemma~\ref{lem:cross-fitting}. For \eqref{eq:CS-remainder}: writing $\delta_a=\hat\tau_a-\tau_a$,
$\delta_b=\hat\tau_b-\tau_b$, $R_n=\langle \delta_a,\hat C\delta_b\rangle_{L^2(\mathbb{P}_0)}$
with $(\hat C\delta_b)(x)=\int\hat C(x,x')\delta_b(x')d\mathbb{P}_0(x')$; Cauchy--Schwarz and the
operator-norm bound give $|R_n|\le\|\hat C\|_{op}\|\delta_a\|\|\delta_b\|$. The smoothness
accounting for the product rate is the quadratic-functional elbow of A5.
\end{proof}

\subsection{Layers L3--L4: moment-vector EIF and composition}

\begin{lemma}[A7: moment-vector EIF and delta-method composition]
\label{lem:moment-eif-A}
Under A0--A8, the moment vector $m$ has influence function $\psi_m$ obtained by stacking:
\begin{enumerate}[(a)]
\item the linear-moment scores
$\psi_{\mathbb{E}_\mu\Delta_a}(O)=\int\bar a(x)\,\mathrm{IF}_{\tau_a}(O;x)\,d\mathbb{P}_0(x)$
for $a\in\{S,Y\}$;
\item the bilinear scores $\chi_{SS},\chi_{YY},\chi_{SY}$ from Lemma~\ref{lem:bilinear-eif-A};
\item for raw second moments, using $\mathbb{E}_\mu[\Delta_a^2]=\Theta_{aa}+(\mathbb{E}_\mu\Delta_a)^2$, the EIF
$\psi_{\mathbb{E}_\mu\Delta_a^2}=\chi_{aa}+2\,\mathbb{E}_\mu[\Delta_a]\,\psi_{\mathbb{E}_\mu\Delta_a}$.
\end{enumerate}
Then, for smooth $\Psi$ with gradient $\nabla\Psi(m)$ (A7 non-degeneracy),
\begin{equation}
\psi_\Theta(O) = \nabla\Psi(m)^\top\,\psi_m(O).
\label{eq:psi-theta}
\end{equation}
The cross-outcome score $\chi_{SY}$ carries
$\mathbb{E}_{\mathbb{P}_0}[\mathrm{IF}_{\tau_S}(O;\cdot)\,\mathrm{IF}_{\tau_Y}(O;\cdot)]$
(the two CATE influence functions are evaluated on the \emph{same unit} $O$), so
$\mathrm{Var}(\psi_\Theta)$ is \textbf{not} block-diagonal across $S$ and $Y$.
\end{lemma}

\begin{proof}[Proof of Lemma~\ref{lem:moment-eif-A}]
(a) is the linear-functional EIF (mean-zero and Riesz representation as in
Lemma~\ref{lem:bilinear-eif-A}, Step~4, with the degenerate weight $\bar a$). (b) is
Lemma~\ref{lem:bilinear-eif-A}. (c) is the chain rule on
$m\mapsto\Theta_{aa}+(\mathbb{E}_\mu\Delta_a)^2$. Finally \eqref{eq:psi-theta} is the functional
delta method (Lemma~\ref{lem:functional-expansion}) applied to $\Psi$ at $m$, whose Hadamard
differentiability and bounded gradient are A7 (Assumptions~\ref{assm:functional},
\ref{assm:lipschitz}, \ref{assm:nondegeneracy}). Non-block-diagonality is immediate because
$\chi_{SY}$ pairs $\mathrm{IF}_{\tau_S}$ and $\mathrm{IF}_{\tau_Y}$ at the same $O$.
\end{proof}

\subsection{Main result: general asymptotic normality}

\begin{theorem}[A: general asymptotic normality, conditional on the geometry]
\label{thm:general}
Under Assumptions~\ref{assm:A0-transport} (A0), \ref{assm:A5-elbow} (A5), and the general-case
readings of Assumptions~\ref{assm:treatment-effects}--\ref{assm:functional} (A1--A4),
\ref{assm:nuisance}--\ref{assm:sample-size} (A6--A8), with one-step cross-fit debiased
estimators of the moments $m$ (Lemmas~\ref{lem:cate-linearity-A}--\ref{lem:elbow}) and
$\hat\Theta=\Psi(\hat m)$, conditional on the sampled geometry $\hat{\mu}_M$ (equivalently on
the kernel $C$), as $n\to\infty$,
\begin{equation}
\sqrt{n}\big(\hat\Theta - \Theta(\hat{\mu}_M)\big)
 = \frac{1}{\sqrt{n}}\sum_{i=1}^n \psi_\Theta(O_i;\hat{\mu}_M)
 + O_P\!\big((n/M)^{1/2}\big) + o_P(1)
 \xrightarrow{d} N\!\big(0,\ \mathbb{E}_{\mathbb{P}_0}[\psi_\Theta^2]\big),
\end{equation}
with $\psi_\Theta$ of \eqref{eq:psi-theta}. If additionally $n=o(M)$
(Assumption~\ref{assm:sample-size}), the MCMC term is $o_P(1)$ and unconditionally
$\sqrt{n}(\hat\Theta-\Theta(\mu))\xrightarrow{d}N(0,\sigma^2(\mu))$,
$\sigma^2(\mu)=\mathbb{E}_{\mathbb{P}_0}[\psi_\Theta(O;\mu)^2]$.
\end{theorem}

\begin{proof}[Proof of Theorem~\ref{thm:general}]
\emph{Step 1 (moment expansions).} For the linear moments, Lemma~\ref{lem:cate-linearity-A}
integrated against $\bar a$ gives
$\widehat{\mathbb{E}_\mu\Delta_a}-\mathbb{E}_\mu\Delta_a=\frac1n\sum_i\psi_{\mathbb{E}_\mu\Delta_a}(O_i)+o_P(n^{-1/2})$.
For the bilinear moments, Lemma~\ref{lem:elbow} gives
$\hat\Theta_{ab}-\Theta_{ab}=\frac1n\sum_i\chi_{ab}(O_i)+o_P(n^{-1/2})$ \emph{provided A5 holds}
(the only use of A5; below the elbow this step fails and the rate degrades). Stacking,
$\hat m-m=\frac1n\sum_i\psi_m(O_i)+o_P(n^{-1/2})$.

\emph{Step 2 (delta method / L4).} $\Psi$ is Hadamard differentiable with bounded gradient on
$\{\mathrm{Var}_\mu\Delta_S,\mathrm{Var}_\mu\Delta_Y\ge c_0\}$ (Assumption~\ref{assm:nondegeneracy});
Lemma~\ref{lem:functional-expansion} (reused, with $m$ in place of the per-study pair) gives
$\sqrt n(\hat\Theta-\Theta(\hat\mu_M)) = \nabla\Psi(m)^\top\sqrt n(\hat m-m)+o_P(1)
= \frac1{\sqrt n}\sum_i\psi_\Theta(O_i;\hat\mu_M)+o_P(1)$ by Lemma~\ref{lem:moment-eif-A}. For the
ratio-type $\Psi$ (correlation), the second-order delta-method remainder is
$O_P(\|\hat m-m\|^2)=O_P(n^{-1})=o_P(n^{-1/2})$ since $\nabla^2\Psi$ is bounded on the
non-degenerate region (A7).
% NOTE(auditor): ratio-Psi second-order remainder closes iff bounded Hessian x O_P(n^{-1}) =
% o_P(n^{-1/2}), i.e. under A7 non-degeneracy. If c_0 -> 0 a TMLE targeting step substitutes for
% one-step + delta method. This is the ratio-sufficiency point to re-check.

\emph{Step 3 (CLT, conditional on $\hat\mu_M$).} Conditional on the geometry, $O_1,\dots,O_n$ are
i.i.d.\ and independent of $\hat\mu_M$; $\mathbb{E}_{\mathbb{P}_0}[\psi_\Theta]=0$ (each stacked
score is mean-zero: Lemma~\ref{lem:bilinear-eif-A}, Step~4) and
$\mathbb{E}_{\mathbb{P}_0}[\psi_\Theta^2]<\infty$ under Assumption~\ref{assm:bounded-if-variance}
(bounded weights, bounded $C$, bounded gradient). The Lindeberg--L\'evy CLT gives
$\frac1{\sqrt n}\sum_i\psi_\Theta(O_i;\hat\mu_M)\xrightarrow{d}N(0,\mathbb{E}[\psi_\Theta^2])$.

\emph{Step 4 (MCMC term).} $\Theta(\hat\mu_M)-\Theta(\mu)=O_P(M^{-1/2})$ by
Lemma~\ref{lem:mcmc-error} applied to the bounded functional $g=\Psi(m(\cdot))$ of the chain
(bounded under A3--A4, A7); scaling by $\sqrt n$ gives the $O_P((n/M)^{1/2})$ term, which is
$o_P(1)$ iff $n=o(M)$ (Assumption~\ref{assm:sample-size}), yielding the unconditional statement.
Combining Steps 1--4 gives the display.
\end{proof}

\subsection{Finite-support special case}

\begin{corollary}[A10: recovery of Theorem~\ref{thm:asymptotic-normality}]
\label{cor:finite-support}
When $X\in\{1,\dots,K\}$ is finite, $C\to\Sigma$ (with $\Sigma_{kk'}=\mathrm{Cov}_\mu(q_k,q_{k'})$),
$\mathrm{IF}_{\tau_a}$ is the per-cell AIPW score, and $\Psi$ is the correlation map $\varphi$ of
Section~\ref{sec:correlation-application}, Theorem~\ref{thm:general} reduces to
Theorem~\ref{thm:asymptotic-normality}, with the same influence function $\psi_\Theta$ and the
same $\sigma^2(\hat\mu_M)$; the $\chi_{ab}$ centering becomes the discrete $\mathrm{tr}(\Sigma V)$
debiasing form, where $V_{kk'}=\mathbb{E}_{\mathbb{P}_0}[\mathrm{IF}_{\tau_a}(O;k)\mathrm{IF}_{\tau_b}(O;k')]$
is the CATE-score covariance.
\end{corollary}

\begin{proof}[Proof of Corollary~\ref{cor:finite-support}]
On finite $X$, A5 is automatic (Remark~\ref{rem:A5-conservative}) and
$\Theta_{ab}=\tau_a^\top\Sigma\tau_b$ by \eqref{eq:discrete-quadform}. The directional representer
satisfies $(h_b(k)p_0(k))_k \propto \Sigma\tau_b$, so Lemma~\ref{lem:bilinear-eif-A} gives the
(mean-zero) score
$\chi_{ab}(O)=(\Sigma\tau_b)^\top\mathbf{IF}_{\tau_a}(O)+(\Sigma\tau_a)^\top\mathbf{IF}_{\tau_b}(O)$,
with $\mathbf{IF}_{\tau_a}(O)=(\mathrm{IF}_{\tau_a}(O;1),\dots,\mathrm{IF}_{\tau_a}(O;K))^\top$
(the plug-in $\tau_a^\top\Sigma\tau_b$ is carried by the plug-in term of \eqref{eq:onestep}).
The population mean of the first-order plug-in bias contracts the CATE-score vectors against
$\Sigma$, giving $\mathrm{tr}(\Sigma V)$ with
$V=\mathbb{E}_{\mathbb{P}_0}[\mathbf{IF}_{\tau_a}\mathbf{IF}_{\tau_b}^\top]$---the disattenuation
term of the finite-support estimator. The per-cell AIPW score is that of
Lemma~\ref{lem:cross-fitting} restricted to a cell, and $\varphi$'s gradient is the composed
gradient of Remark~\ref{rem:aggregate-gradient}; substituting into \eqref{eq:psi-theta} reproduces
$\psi_\Theta(O;\hat\mu_M)$ of Theorem~\ref{thm:asymptotic-normality}. All remainders are those of
Lemmas~\ref{lem:cross-fitting}, \ref{lem:functional-expansion}, so the two limit laws coincide.
% NOTE(auditor): The exact tr(Sigma V) constant depends on the factor-2 diagonal vs off-diagonal
% pairing (Lemma \ref{lem:bilinear-eif-A}, Step 5). The claim is STRUCTURAL collapse (same object),
% not a re-derivation of the finite-support constant (already validated in
% Theorem \ref{thm:asymptotic-normality}). Re-check pairing if a numeric match is asserted.
\end{proof}

\subsection{Post-proof audit (general theory)}
\label{subsec:audit-general}

\paragraph{Assumption check.}
A0 (Assumption~\ref{assm:A0-transport}) is used once, to license the frozen-effect
representation \eqref{eq:delta-linear}. A1--A2 enter Lemma~\ref{lem:cate-linearity-A}; A3
(Assumption~\ref{assm:treatment-effects}(c)) bounds $C$ and the weights; A4
(Assumption~\ref{assm:geometry}) gives the bounded kernel and MCMC-CLT; A5
(Assumption~\ref{assm:A5-elbow}) is used once, in Lemma~\ref{lem:elbow}; A6 controls nuisance
remainders; A7 (Assumptions~\ref{assm:functional}, \ref{assm:lipschitz}, \ref{assm:nondegeneracy})
gives the delta method and bounded Hessian; A8 (Assumption~\ref{assm:sample-size}) governs the $n$
vs $M$ regime. No assumption is introduced mid-proof.

\paragraph{Dependency check.}
Riesz representation (Lemma~\ref{lem:bilinear-eif-A}): square-integrable CATE scores and bounded
$C$, verified. Fubini/Leibniz: bounded $C$, $L^2$ CATEs, verified. Cauchy--Schwarz with
operator-norm bound (Lemma~\ref{lem:elbow}): symmetric bounded kernel, verified under A3--A4.
Chernozhukov et al.\ (2018) orthogonality/cross-fitting: reused via Lemmas~\ref{lem:orthogonality},
\ref{lem:cross-fitting} with the $\eta$-free $h_a,h_b$ in the role of $w$. Markov-chain CLT
(Lemma~\ref{lem:mcmc-error}) and delta method (Lemma~\ref{lem:functional-expansion}) reused
unchanged. Minimax lower bounds
\cite{BickelRitov1988,BirgeMassart1995,Robins2008HOIF,Robins2017Minimax,Kennedy2024Minimax,McClean2024DCDR}
supply the below-elbow non-existence claim.

\paragraph{Rate verification.}
Bilinear remainder \eqref{eq:CS-remainder}: product of CATE $L^2$-rates, $o_P(n^{-1/2})$ iff
$s_S+s_Y>d/2$ (A5); at equality $O_P(n^{-1/2})$ (no slack), below strictly slower and dominant, so
no $\sqrt n$-normal limit. Nuisance-$\eta$ remainder: $o_P(n^{-1/2})$ via A6, as in
Lemma~\ref{lem:cross-fitting}. Delta-method second-order remainder: $O_P(n^{-1})=o_P(n^{-1/2})$
under bounded Hessian (A7). MCMC: $O_P((n/M)^{1/2})$, $o_P(1)$ iff $n=o(M)$.

\paragraph{Edge-case scan.}
(i) \emph{Near-elbow} $s_S+s_Y\downarrow d/2$: product rate loses slack, intervals widen; report
the minimax rate honestly. (ii) \emph{Weak overlap} $e(x)\to0,1$: $\mathrm{IF}_{\tau_a}$ blows up
(Assumption~\ref{assm:second-derivatives} fails); trim. (iii) \emph{Degenerate variance}: $\nabla\Psi$
blows up like $c_0^{-3/2}$; A7 excludes. (iv) \emph{Kernel conditioning}: the bound
\eqref{eq:CS-remainder} uses $\|C\|_{op}$, not $C^{-1}$, so ill-conditioning of $C$ does not break
it---only extreme weights (bounded by A3) matter.

\paragraph{Tightness.}
Above the elbow the $\sqrt n$ rate is sharp: $\psi_\Theta$ is the canonical gradient, so
$\mathbb{E}[\psi_\Theta^2]$ is the semiparametric efficiency bound. Below the elbow the
$n^{-4s/(4s+d)}$ rate is minimax, so A5 is not improvable as a $\sqrt n$-condition for the Dirac
kernel; for the smooth kernel $C$ it is conservative (Remark~\ref{rem:A5-conservative}).

\paragraph{Failure localization.}
Three steps warrant re-checking. (1) The $a\ne b$ symmetrization
(Lemma~\ref{lem:bilinear-eif-A}, Steps~2--4): the two directional representers must be added,
guarded by the diagonal (Step~5) and discrete (Step~6) checks. (2) Kernel/covariate-law
$\mathbb{P}_0$-dependence (Remark~\ref{rem:kernel-dependence}): explicitly deferred. (3) Ratio
one-step sufficiency (Theorem~\ref{thm:general}, Step~2): closes under bounded Hessian (A7); TMLE
fallback if $c_0\to0$. The linear/bilinear/CLT core is otherwise standard.


\bibliography{refs}
\bibliographystyle{plain}
\end{document}
